% Chapter 9. All quoted numbers come from chapters/ch09-numbers.tex, written by code/ch09/run_all.py.
% Generated by code/ch09/run_all.py. Do not edit by hand.
\newcommand{\bfAfterMean}{7.2}
\newcommand{\bfEnterShare}{82}
\newcommand{\bfIncMean}{9.6}
\newcommand{\bfIncTheory}{9.4}
\newcommand{\bfRecordBias}{0.7}
\newcommand{\bfRecordTheory}{0.6}
\newcommand{\btArith}{71}
\newcommand{\btBest}{66}
\newcommand{\btBetaLate}{2.0}
\newcommand{\btCorrEarly}{0.10}
\newcommand{\btCorrEth}{0.60}
\newcommand{\btCorrLate}{0.55}
\newcommand{\btCovSign}{0.33}
\newcommand{\btFirst}{January 2015}
\newcommand{\btFirstPrice}{December 2014}
\newcommand{\btGeo}{60}
\newcommand{\btKurt}{0.4}
\newcommand{\btLast}{September 2026}
\newcommand{\btMddDaily}{84}
\newcommand{\btMeanSe}{21}
\newcommand{\btMonths}{141}
\newcommand{\btNextDay}{8 January 2015}
\newcommand{\btPriceEnd}{83,482}
\newcommand{\btPriceStart}{340}
\newcommand{\btSR}{0.97}
\newcommand{\btSeriesStart}{1 December 2014}
\newcommand{\btShareFive}{20}
\newcommand{\btShareOne}{3}
\newcommand{\btShareTwo}{6}
\newcommand{\btSkew}{0.55}
\newcommand{\btStartDay}{18 December 2014}
\newcommand{\btUnderYears}{3.0}
\newcommand{\btVol}{72}
\newcommand{\btVolFirstYear}{115}
\newcommand{\btVolFirstYearName}{2015}
\newcommand{\btVolLastYear}{42}
\newcommand{\btVolLastYearName}{2025}
\newcommand{\btVolSixty}{9.9}
\newcommand{\btWFive}{1.6}
\newcommand{\btWHalf}{12.1}
\newcommand{\btWQuarter}{6.2}
\newcommand{\btWTen}{2.9}
\newcommand{\btWorst}{\ensuremath{-}38}
\newcommand{\btYearsSig}{4}
\newcommand{\cadCarry}{0.04}
\newcommand{\cadCarryMax}{1.8}
\newcommand{\cadCarryMin}{\ensuremath{-}1.8}
\newcommand{\cadCarryPos}{62}
\newcommand{\cadCross}{0.00}
\newcommand{\cadEx}{\ensuremath{-}1.04}
\newcommand{\cadFirst}{June 2007}
\newcommand{\cadIdGap}{\ensuremath{4\times 10^{-16}}}
\newcommand{\cadLast}{September 2026}
\newcommand{\cadMonths}{232}
\newcommand{\cadSlope}{\ensuremath{-}2.3}
\newcommand{\cadSlopeSe}{2.6}
\newcommand{\cadSpot}{\ensuremath{-}1.08}
\newcommand{\cadVol}{8.7}
\newcommand{\ccAMdd}{22}
\newcommand{\ccAMean}{6.9}
\newcommand{\ccASR}{0.69}
\newcommand{\ccATwo}{\ensuremath{-}18}
\newcommand{\ccATwoAbs}{18}
\newcommand{\ccAVol}{9.9}
\newcommand{\ccAWorst}{\ensuremath{-}7.8}
\newcommand{\ccBtcTwo}{\ensuremath{-}64}
\newcommand{\ccBtcTwoAbs}{64}
\newcommand{\ccFiveDiff}{0.20}
\newcommand{\ccFiveHi}{0.42}
\newcommand{\ccFiveLo}{0.03}
\newcommand{\ccFiveMdd}{24}
\newcommand{\ccFiveMean}{10.0}
\newcommand{\ccFiveSR}{0.90}
\newcommand{\ccFiveVol}{11.1}
\newcommand{\ccFiveWorst}{\ensuremath{-}8.0}
\newcommand{\ccTwoDiff}{0.10}
\newcommand{\ccTwoHi}{0.21}
\newcommand{\ccTwoLo}{0.02}
\newcommand{\ccTwoMdd}{23}
\newcommand{\ccTwoMean}{8.1}
\newcommand{\ccTwoSR}{0.79}
\newcommand{\ccTwoTwo}{\ensuremath{-}19}
\newcommand{\ccTwoTwoAbs}{19}
\newcommand{\ccTwoVol}{10.3}
\newcommand{\ccTwoWorst}{\ensuremath{-}7.7}
\newcommand{\ceArith}{110}
\newcommand{\ceBest}{206}
\newcommand{\ceFirst}{June 2016}
\newcommand{\ceGeo}{66}
\newcommand{\ceKurt}{8.7}
\newcommand{\ceMddDaily}{94}
\newcommand{\ceMeanSe}{40}
\newcommand{\ceMonths}{124}
\newcommand{\ceSR}{0.84}
\newcommand{\ceSkew}{2.26}
\newcommand{\ceUnderYears}{3.8}
\newcommand{\ceVol}{128}
\newcommand{\ceWorst}{\ensuremath{-}54}
\newcommand{\cfChangeA}{0.12}
\newcommand{\cfChangeB}{0.03}
\newcommand{\cfFlipsMed}{25}
\newcommand{\cfMonths}{284}
\newcommand{\cfRangeA}{6.6}
\newcommand{\cfRatio}{3.9}
\newcommand{\cfTwelveMax}{1.07}
\newcommand{\cfTwelveMin}{0.45}
\newcommand{\clB}{0.511}
\newcommand{\clDrift}{0}
\newcommand{\clFunds}{20,000}
\newcommand{\clLoss}{40}
\newcommand{\clMaxGapBgk}{0.1}
\newcommand{\clMeanBgkFive}{1.33}
\newcommand{\clMeanBgkTen}{1.87}
\newcommand{\clMeanFive}{1.50}
\newcommand{\clMeanSimFive}{1.34}
\newcommand{\clMeanSimTen}{1.91}
\newcommand{\clMeanTen}{2.00}
\newcommand{\clSig}{15}
\newcommand{\clSurvBgkFive}{89.0}
\newcommand{\clSurvBgkTen}{74.2}
\newcommand{\clSurvFive}{87.2}
\newcommand{\clSurvSe}{0.3}
\newcommand{\clSurvSimFive}{89.0}
\newcommand{\clSurvSimTen}{74.1}
\newcommand{\clSurvTen}{71.8}
\newcommand{\clYears}{10}
\newcommand{\csAGfc}{\ensuremath{-}29}
\newcommand{\csAGfcAbs}{29}
\newcommand{\csAMdd}{30}
\newcommand{\csAMean}{5.8}
\newcommand{\csASR}{0.63}
\newcommand{\csATwo}{\ensuremath{-}18}
\newcommand{\csATwoAbs}{18}
\newcommand{\csAVol}{9.2}
\newcommand{\csAWorst}{\ensuremath{-}10.7}
\newcommand{\csBGfc}{\ensuremath{-}24}
\newcommand{\csBGfcAbs}{24}
\newcommand{\csBMdd}{26}
\newcommand{\csBMean}{5.3}
\newcommand{\csBSR}{0.65}
\newcommand{\csBTwo}{\ensuremath{-}14}
\newcommand{\csBTwoAbs}{14}
\newcommand{\csBVol}{8.2}
\newcommand{\csBWorst}{\ensuremath{-}8.7}
\newcommand{\csBlock}{12}
\newcommand{\csCorrAT}{\ensuremath{-}0.19}
\newcommand{\csDiff}{0.026}
\newcommand{\csDiffG}{0.06}
\newcommand{\csDiffGHi}{0.12}
\newcommand{\csDiffGLo}{\ensuremath{-}0.01}
\newcommand{\csDiffHi}{0.08}
\newcommand{\csDiffLo}{\ensuremath{-}0.04}
\newcommand{\csDiffP}{21}
\newcommand{\csDiffTab}{0.03}
\newcommand{\csGMdd}{25}
\newcommand{\csGMean}{5.6}
\newcommand{\csGSR}{0.69}
\newcommand{\csGVol}{8.2}
\newcommand{\csGWorst}{\ensuremath{-}8.7}
\newcommand{\csReps}{5,000}
\newcommand{\cwEnd}{\ensuremath{-}9.2}
\newcommand{\cwEndAbs}{9.2}
\newcommand{\cwEndGross}{\ensuremath{-}2.2}
\newcommand{\cwEndGrossAbs}{2.2}
\newcommand{\cwFeeWedge}{7.0}
\newcommand{\cwFromPeak}{16.8}
\newcommand{\cwFromPeakUp}{17}
\newcommand{\cwGfc}{6.0}
\newcommand{\cwPeak}{7.5}
\newcommand{\cwPeakWhen}{February 2009}
\newcommand{\dbBack}{8.18}
\newcommand{\dbBackBias}{0.54}
\newcommand{\dbClean}{7.63}
\newcommand{\dbCleanBias}{0.28}
\newcommand{\dbDeadShare}{74}
\newcommand{\dbEventLast}{36}
\newcommand{\dbFunds}{3,000}
\newcommand{\dbInc}{24}
\newcommand{\dbLive}{8.91}
\newcommand{\dbMonths}{240}
\newcommand{\dbNall}{2,880}
\newcommand{\dbNdb}{2,315}
\newcommand{\dbNlive}{594}
\newcommand{\dbPostEntry}{7.40}
\newcommand{\dbPreEntry}{9.90}
\newcommand{\dbSurvBias}{1.56}
\newcommand{\dbTruth}{7.35}
\newcommand{\dcAbsCad}{0.18}
\newcommand{\dcAbsEq}{0.06}
\newcommand{\dcAbsEur}{0.12}
\newcommand{\dcAbsGoc}{\ensuremath{-}0.15}
\newcommand{\dcAbsJpy}{0.06}
\newcommand{\dcAbsOil}{\ensuremath{-}0.02}
\newcommand{\dcAbsPos}{5}
\newcommand{\dcAbsSig}{0}
\newcommand{\dcAbsTCad}{1.4}
\newcommand{\dcAbsTEq}{0.8}
\newcommand{\dcAbsTEur}{1.0}
\newcommand{\dcAbsTGoc}{\ensuremath{-}0.7}
\newcommand{\dcAbsTJpy}{0.5}
\newcommand{\dcAbsTOil}{\ensuremath{-}0.7}
\newcommand{\dcAbsTUst}{0.5}
\newcommand{\dcAbsUst}{0.09}
\newcommand{\dcAlphaA}{5.2}
\newcommand{\dcAlphaB}{\ensuremath{-}3.4}
\newcommand{\dcBetaCad}{0.01}
\newcommand{\dcBetaEq}{\ensuremath{-}0.14}
\newcommand{\dcBetaEur}{0.13}
\newcommand{\dcBetaGoc}{\ensuremath{-}0.37}
\newcommand{\dcBetaJpy}{\ensuremath{-}0.21}
\newcommand{\dcBetaOil}{\ensuremath{-}0.06}
\newcommand{\dcBetaUst}{0.54}
\newcommand{\dcIdGap}{\ensuremath{1\times 10^{-16}}}
\newcommand{\dcLoadPart}{7.4}
\newcommand{\dcLoadPartA}{\ensuremath{-}1.2}
\newcommand{\dcLoadShift}{8.6}
\newcommand{\dcMean}{4.00}
\newcommand{\dcMonths}{308}
\newcommand{\dcRegMonths}{307}
\newcommand{\dcRtwoA}{0.15}
\newcommand{\dcRtwoB}{0.17}
\newcommand{\dcSlCad}{0.42}
\newcommand{\dcSlEq}{1.36}
\newcommand{\dcSlEur}{\ensuremath{-}0.03}
\newcommand{\dcSlGoc}{0.56}
\newcommand{\dcSlJpy}{1.15}
\newcommand{\dcSlOil}{\ensuremath{-}0.14}
\newcommand{\dcSlUst}{0.67}
\newcommand{\dcSumGap}{0.0}
\newcommand{\dcTA}{2.4}
\newcommand{\dcTB}{\ensuremath{-}0.7}
\newcommand{\dqBottom}{0.13}
\newcommand{\dqMid}{0.06}
\newcommand{\dqObs}{2,143}
\newcommand{\dqRho}{0.008}
\newcommand{\dqTailShare}{\ensuremath{-}28}
\newcommand{\dqTailShareAbs}{28}
\newcommand{\dqTop}{\ensuremath{-}0.01}
\newcommand{\exAPnl}{1.0}
\newcommand{\exAQ}{34}
\newcommand{\exAS}{6}
\newcommand{\exBMonth}{0.16}
\newcommand{\exBSR}{0.14}
\newcommand{\exCLim}{0.63}
\newcommand{\exCSR}{0.53}
\newcommand{\exDBias}{4.4}
\newcommand{\exDMean}{10.4}
\newcommand{\exDShare}{73}
\newcommand{\exEBias}{0.98}
\newcommand{\exEInc}{10.9}
\newcommand{\exFInf}{1.40}
\newcommand{\exFT}{111}
\newcommand{\exGFee}{3}
\newcommand{\exGNet}{2}
\newcommand{\exGRate}{30}
\newcommand{\exHShare}{10}
\newcommand{\exHSharex}{10.34}
\newcommand{\exHW}{6.1}
\newcommand{\exHWx}{6.085}
\newcommand{\exINet}{5.6}
\newcommand{\exINetEx}{2.6}
\newcommand{\exIRatio}{43}
\newcommand{\exJBest}{59}
\newcommand{\exJMax}{1.87}
\newcommand{\exKOne}{5.0}
\newcommand{\exKReset}{10.76}
\newcommand{\exKResetThree}{5.76}
\newcommand{\exKRise}{23.84}
\newcommand{\exKThree}{0.96}
\newcommand{\exKTot}{5.96}
\newcommand{\exLHSix}{0.2}
\newcommand{\exLHTen}{1.0}
\newcommand{\exLNSix}{1.2}
\newcommand{\exLNTen}{2.0}
\newcommand{\exLSSix}{1.0}
\newcommand{\exLSTen}{2.0}
\newcommand{\exMCarry}{\ensuremath{-}1.5}
\newcommand{\exMCross}{0.045}
\newcommand{\exMEx}{0.049}
\newcommand{\exMSpot}{1.504}
\newcommand{\exNNak}{0.0143}
\newcommand{\exNRuin}{0.000244}
\newcommand{\feCatch}{6.25}
\newcommand{\feExFeeH}{4.00}
\newcommand{\feExFeeR}{9.22}
\newcommand{\feExMark}{16.00}
\newcommand{\feExNavH}{13.10}
\newcommand{\feExNavR}{7.88}
\newcommand{\feF}{20}
\newcommand{\feH}{5}
\newcommand{\feHard}{1.29}
\newcommand{\feHardSim}{1.29}
\newcommand{\feIdGap}{\ensuremath{9\times 10^{-16}}}
\newcommand{\feMu}{8}
\newcommand{\feNone}{1.96}
\newcommand{\feNoneSim}{1.96}
\newcommand{\feOneCall}{1.50}
\newcommand{\feOptApprox}{0.80}
\newcommand{\feOptAtm}{1.12}
\newcommand{\feOptFwd}{0.80}
\newcommand{\feOptVol}{10}
\newcommand{\feRf}{3}
\newcommand{\feSig}{12}
\newcommand{\feSimSe}{0.002}
\newcommand{\feSims}{40,000}
\newcommand{\feSoft}{1.87}
\newcommand{\feSoftSim}{1.87}
\newcommand{\feTenClosed}{14.0}
\newcommand{\feTenHwm}{9.5}
\newcommand{\feTenRatio}{68}
\newcommand{\feTenReset}{14.0}
\newcommand{\feTenResetSe}{0.05}
\newcommand{\feTenVol}{15}
\newcommand{\feYears}{10}
\newcommand{\ffAllFees}{3.78}
\newcommand{\ffFeeShare}{47}
\newcommand{\ffFmg}{1}
\newcommand{\ffFpf}{10}
\newcommand{\ffLossYears}{19}
\newcommand{\ffLossYearsPaid}{18}
\newcommand{\ffMu}{8}
\newcommand{\ffNetFof}{4.22}
\newcommand{\ffNetOne}{5.84}
\newcommand{\ffPerf}{20}
\newcommand{\ffPoolFive}{1.84}
\newcommand{\ffPoolLimit}{1.74}
\newcommand{\ffPoolOne}{2.16}
\newcommand{\ffPoolTen}{1.79}
\newcommand{\ffPoolTwenty}{1.77}
\newcommand{\ffRateLimit}{22}
\newcommand{\ffRateSep}{27.0}
\newcommand{\ffRho}{0.3}
\newcommand{\ffSepFive}{2.16}
\newcommand{\ffSepOne}{2.16}
\newcommand{\ffSepTen}{2.16}
\newcommand{\ffSepTwenty}{2.16}
\newcommand{\ffSig}{15}
\newcommand{\ffSimPool}{1.79}
\newcommand{\ffSimSep}{2.16}
\newcommand{\ffSimSepSe}{0.003}
\newcommand{\ffVolFive}{9.9}
\newcommand{\ffVolLimit}{8.2}
\newcommand{\ffVolOne}{15.0}
\newcommand{\ffVolTen}{9.1}
\newcommand{\ffVolTwenty}{8.7}
\newcommand{\gcCarry}{1.05}
\newcommand{\gcCarryShare}{44}
\newcommand{\gcExp}{1.77}
\newcommand{\gcExpPos}{87}
\newcommand{\gcExpShare}{74}
\newcommand{\gcFirst}{February 2001}
\newcommand{\gcLast}{August 2026}
\newcommand{\gcReal}{2.39}
\newcommand{\gcRest}{0.62}
\newcommand{\gcRoll}{0.72}
\newcommand{\gcSlope}{1.7}
\newcommand{\gcSlopeSe}{0.8}
\newcommand{\gcVolExp}{0.44}
\newcommand{\gcVolReal}{6.2}
\newcommand{\gcYEnd}{3.73}
\newcommand{\gcYStart}{5.39}
\newcommand{\hdActualDrag}{3.1}
\newcommand{\hdAlphaQ}{1.44}
\newcommand{\hdAlphaT}{2.0}
\newcommand{\hdBBd}{0.008}
\newcommand{\hdBCom}{0.026}
\newcommand{\hdBFx}{\ensuremath{-}0.042}
\newcommand{\hdBIr}{0.018}
\newcommand{\hdBStk}{\ensuremath{-}0.015}
\newcommand{\hdCanGate}{45.1}
\newcommand{\hdCanGateFirst}{51.9}
\newcommand{\hdCountFirst}{1,219}
\newcommand{\hdCountLast}{2,257}
\newcommand{\hdCreditCovid}{7.3}
\newcommand{\hdCushionMin}{36.1}
\newcommand{\hdCushionMinQ}{2017Q4}
\newcommand{\hdDragCredit}{3.1}
\newcommand{\hdDragEquity}{2.7}
\newcommand{\hdDragEvent}{6.0}
\newcommand{\hdDragFutures}{2.3}
\newcommand{\hdDragMacro}{2.1}
\newcommand{\hdDragMulti}{2.8}
\newcommand{\hdDragOther}{6.1}
\newcommand{\hdDragRv}{2.4}
\newcommand{\hdFeeInc}{11.5}
\newcommand{\hdFeeIncSe}{1.0}
\newcommand{\hdFeeMgmt}{1.41}
\newcommand{\hdFeeMgmtSe}{0.13}
\newcommand{\hdFeeNegDrag}{1.48}
\newcommand{\hdFeeNegObs}{52}
\newcommand{\hdFeeNegQ}{52}
\newcommand{\hdFeeObs}{376}
\newcommand{\hdFeeRsq}{0.38}
\newcommand{\hdFirst}{2013Q1}
\newcommand{\hdGatedLast}{2.4}
\newcommand{\hdGatedMax}{2.6}
\newcommand{\hdGatedMaxQ}{2025Q1}
\newcommand{\hdGavLast}{13.9}
\newcommand{\hdInvLong}{64.4}
\newcommand{\hdInvNinety}{35.7}
\newcommand{\hdInvNinetyMean}{42.4}
\newcommand{\hdInvSeven}{8.4}
\newcommand{\hdInvSevenMean}{11.3}
\newcommand{\hdLast}{2026Q2}
\newcommand{\hdLevAgg}{2.5}
\newcommand{\hdLevCredit}{1.7}
\newcommand{\hdLevEquity}{1.8}
\newcommand{\hdLevEvent}{1.3}
\newcommand{\hdLevFutures}{2.5}
\newcommand{\hdLevMacro}{6.8}
\newcommand{\hdLevMax}{8.1}
\newcommand{\hdLevMaxName}{relative value}
\newcommand{\hdLevMulti}{4.5}
\newcommand{\hdLevOther}{1.4}
\newcommand{\hdLevRv}{8.1}
\newcommand{\hdMeanCredit}{11.5}
\newcommand{\hdMeanEquity}{13.2}
\newcommand{\hdMeanEvent}{13.6}
\newcommand{\hdMeanFutures}{6.0}
\newcommand{\hdMeanMacro}{5.9}
\newcommand{\hdMeanMulti}{9.2}
\newcommand{\hdMeanOther}{18.6}
\newcommand{\hdMeanRv}{7.5}
\newcommand{\hdNCredit}{54}
\newcommand{\hdNEquity}{54}
\newcommand{\hdNEvent}{54}
\newcommand{\hdNFutures}{52}
\newcommand{\hdNMacro}{54}
\newcommand{\hdNMulti}{54}
\newcommand{\hdNOther}{54}
\newcommand{\hdNRv}{54}
\newcommand{\hdNavFirst}{2.0}
\newcommand{\hdNavLast}{5.5}
\newcommand{\hdPorLong}{19.3}
\newcommand{\hdPorNinety}{80.7}
\newcommand{\hdPorNinetyMean}{84.1}
\newcommand{\hdPorOne}{41.2}
\newcommand{\hdPorSeven}{64.8}
\newcommand{\hdPorSevenMean}{63.6}
\newcommand{\hdProgCorr}{0.37}
\newcommand{\hdProgN}{52}
\newcommand{\hdPtfsAdj}{\ensuremath{-}0.01}
\newcommand{\hdPtfsCovidMin}{99}
\newcommand{\hdPtfsLast}{August 2026}
\newcommand{\hdPtfsN}{52}
\newcommand{\hdQuarters}{54}
\newcommand{\hdRhoCredit}{0.12}
\newcommand{\hdRhoEquity}{\ensuremath{-}0.08}
\newcommand{\hdRhoEvent}{0.26}
\newcommand{\hdRhoFutures}{0.07}
\newcommand{\hdRhoMacro}{0.07}
\newcommand{\hdRhoMulti}{0.02}
\newcommand{\hdRhoOther}{0.29}
\newcommand{\hdRhoRv}{\ensuremath{-}0.07}
\newcommand{\hdRhoSe}{0.14}
\newcommand{\hdRsqCredit}{0.43}
\newcommand{\hdRsqCreditEx}{0.10}
\newcommand{\hdRsqEquity}{0.37}
\newcommand{\hdRsqEvent}{0.37}
\newcommand{\hdRsqFutures}{0.09}
\newcommand{\hdRsqMacro}{0.09}
\newcommand{\hdRsqMulti}{0.29}
\newcommand{\hdRsqOther}{0.39}
\newcommand{\hdRsqRv}{0.54}
\newcommand{\hdRvCovid}{9.0}
\newcommand{\hdShareCredit}{21}
\newcommand{\hdShareEquity}{17}
\newcommand{\hdShareEvent}{31}
\newcommand{\hdShareFutures}{28}
\newcommand{\hdShareMacro}{26}
\newcommand{\hdShareMulti}{24}
\newcommand{\hdShareOther}{25}
\newcommand{\hdShareRv}{24}
\newcommand{\hdTBd}{0.5}
\newcommand{\hdTCom}{0.9}
\newcommand{\hdTFx}{\ensuremath{-}1.6}
\newcommand{\hdTIr}{1.1}
\newcommand{\hdTStk}{\ensuremath{-}0.7}
\newcommand{\hdTwoTwentyDrag}{4.5}
\newcommand{\hdVolCredit}{4.5}
\newcommand{\hdVolEquity}{11.1}
\newcommand{\hdVolEvent}{7.6}
\newcommand{\hdVolFutures}{9.5}
\newcommand{\hdVolMacro}{4.3}
\newcommand{\hdVolMulti}{4.7}
\newcommand{\hdVolOther}{5.5}
\newcommand{\hdVolRv}{4.1}
\newcommand{\lbSE}{0.20}
\newcommand{\lbSREighteen}{0.31}
\newcommand{\lbSRNine}{0.41}
\newcommand{\lbSROne}{0.05}
\newcommand{\lbSRSix}{0.19}
\newcommand{\lbSRThree}{0.24}
\newcommand{\lbSRTwelve}{0.35}
\newcommand{\lbSRTwentyFour}{0.11}
\newcommand{\ocEarlyHi}{0.34}
\newcommand{\ocEarlyLo}{\ensuremath{-}0.16}
\newcommand{\ocEpochYears}{3.99}
\newcommand{\ocFisherP}{0.003}
\newcommand{\ocFisherZ}{2.95}
\newcommand{\ocHalvings}{33}
\newcommand{\ocLateHi}{0.69}
\newcommand{\ocLateLo}{0.37}
\newcommand{\ocNEarly}{60}
\newcommand{\ocNLate}{80}
\newcommand{\ocNakSix}{0.00024}
\newcommand{\ocPaperGap}{0}
\newcommand{\ocQLast}{0.45}
\newcommand{\ocRatioMax}{12.9}
\newcommand{\ocRatioMin}{2.2}
\newcommand{\ocRatioRecentMax}{4.0}
\newcommand{\ocRatioRecentMaxYear}{2024}
\newcommand{\ocRatioRecentMin}{2.2}
\newcommand{\ocRuinSix}{\ensuremath{1.9\times 10^{-6}}}
\newcommand{\ocRuinTen}{\ensuremath{2.9\times 10^{-10}}}
\newcommand{\ocShareFive}{96.88}
\newcommand{\ocShareFour}{93.75}
\newcommand{\ocShareOne}{50}
\newcommand{\ocShortfall}{0.0231}
\newcommand{\ocSupply}{20,999,999.9769}
\newcommand{\ocVolMax}{115}
\newcommand{\ocVolMaxYear}{2015}
\newcommand{\ocVolMin}{42}
\newcommand{\ocVolMinYear}{2025}
\newcommand{\ocZCount}{8}
\newcommand{\ocZLast}{340}
\newcommand{\ocZMatch}{8}
\newcommand{\ocZThirty}{24}
\newcommand{\oilFOne}{70.7}
\newcommand{\oilRoll}{\ensuremath{-}11.9}
\newcommand{\oilS}{70}
\newcommand{\pmCount}{8}
\newcommand{\pmHeld}{24}
\newcommand{\pmHeldFirst}{February 1999}
\newcommand{\pmHeldLast}{January 2001}
\newcommand{\pmLive}{307}
\newcommand{\pmMean}{0.01}
\newcommand{\pmMeanGross}{0.05}
\newcommand{\pmN}{400}
\newcommand{\pmP}{0.022}
\newcommand{\pmQ}{0.30}
\newcommand{\pmSd}{0.19}
\newcommand{\psEqAprTwenty}{\ensuremath{-}0.15}
\newcommand{\psEqFebTwenty}{0.19}
\newcommand{\psEqMarTwenty}{0.18}
\newcommand{\psGfcFirstShort}{February 2008}
\newcommand{\psGfcMonths}{16}
\newcommand{\psGfcShortEq}{13}
\newcommand{\psMax}{0.74}
\newcommand{\psTwoTwoShortGoc}{12}
\newcommand{\psTwoTwoShortUst}{12}
\newcommand{\pvInSample}{0.708}
\newcommand{\pvK}{7}
\newcommand{\pvSim}{1.429}
\newcommand{\pvSimSe}{0.015}
\newcommand{\pvSims}{20,000}
\newcommand{\pvT}{24}
\newcommand{\pvTheory}{1.438}
\newcommand{\rpCorrA}{0.27}
\newcommand{\rpCorrB}{0.86}
\newcommand{\rpFullA}{0.12}
\newcommand{\rpIdio}{3}
\newcommand{\rpInA}{0.54}
\newcommand{\rpInB}{0.80}
\newcommand{\rpMeanGapA}{5.8}
\newcommand{\rpMeanGapB}{0.4}
\newcommand{\rpMonths}{284}
\newcommand{\rpOutA}{\ensuremath{-}0.31}
\newcommand{\rpOutB}{0.73}
\newcommand{\rpSRa}{\ensuremath{-}0.26}
\newcommand{\rpSRb}{0.28}
\newcommand{\rpSRt}{0.31}
\newcommand{\rpTeA}{12.0}
\newcommand{\rpTeB}{5.5}
\newcommand{\rpWin}{24}
\newcommand{\rsCorrEarly}{0.11}
\newcommand{\rsCorrLate}{0.50}
\newcommand{\rsCross}{1.3}
\newcommand{\rsVolBtcEarly}{81}
\newcommand{\rsVolBtcLate}{64}
\newcommand{\rsVolSixtyEarly}{6.9}
\newcommand{\rsVolSixtyLate}{11.7}
\newcommand{\rsWHalfEarly}{7.8}
\newcommand{\rsWHalfLate}{15.5}
\newcommand{\rsWQuarterEarly}{4.4}
\newcommand{\rsWQuarterLate}{7.4}
\newcommand{\rsWTenEarly}{2.4}
\newcommand{\rsWTenLate}{3.2}
\newcommand{\scBonds}{0.68}
\newcommand{\scMax}{0.68}
\newcommand{\scMaxPair}{US 10 year bond and Canada 10 year bond}
\newcommand{\scMin}{\ensuremath{-}0.09}
\newcommand{\scMinPair}{US equities and US 10 year bond}
\newcommand{\scMonths}{295}
\newcommand{\scNeg}{3}
\newcommand{\scPairs}{21}
\newcommand{\scSdPairs}{0.16}
\newcommand{\slHighP}{0.2}
\newcommand{\slIcpt}{0.56}
\newcommand{\slLin}{\ensuremath{-}0.24}
\newcommand{\slLinT}{\ensuremath{-}3.4}
\newcommand{\slLowP}{4.4}
\newcommand{\slMidP}{0.1}
\newcommand{\slQuad}{1.28}
\newcommand{\slQuadT}{2.2}
\newcommand{\slQuarters}{102}
\newcommand{\srDriftAvg}{0.06}
\newcommand{\srDriftEq}{0.27}
\newcommand{\srPredSRavg}{0.00}
\newcommand{\srRealSRavg}{0.17}
\newcommand{\srRhoAvg}{\ensuremath{-}0.001}
\newcommand{\srRhoEq}{0.046}
\newcommand{\srRhoJpy}{0.058}
\newcommand{\srRhoOil}{\ensuremath{-}0.074}
\newcommand{\srTimingAvg}{0.11}
\newcommand{\srTimingEq}{0.17}
\newcommand{\stLossFinite}{68.17}
\newcommand{\stLossSe}{0.75}
\newcommand{\stLossShare}{66.3}
\newcommand{\stLossTheory}{68.3}
\newcommand{\stLossZ}{2.4}
\newcommand{\stMaxErr}{\ensuremath{2\times 10^{-15}}}
\newcommand{\stMeanPnl}{0.015}
\newcommand{\stNday}{252}
\newcommand{\stNpath}{4,000}
\newcommand{\stSkew}{2.5}
\newcommand{\svBiasZero}{1.3}
\newcommand{\svBiasZeroTh}{1.3}
\newcommand{\svMaxGap}{0.4}
\newcommand{\svMu}{0.6}
\newcommand{\svShareZero}{88}
\newcommand{\svSig}{3}
\newcommand{\svT}{36}
\newcommand{\sxAbove}{1}
\newcommand{\sxActualFine}{0.096}
\newcommand{\sxCount}{15}
\newcommand{\sxMedian}{0.03}
\newcommand{\sxMedianOne}{\ensuremath{-}0.038}
\newcommand{\sxNeeded}{2,200}
\newcommand{\sxQ}{0.07}
\newcommand{\sxReps}{2,000}
\newcommand{\sxSingle}{400,000}
\newcommand{\sxTokens}{20}
\newcommand{\tfBestYear}{34}
\newcommand{\tfBestYearLab}{2008}
\newcommand{\tfBetaEq}{\ensuremath{-}0.18}
\newcommand{\tfCorrEq}{\ensuremath{-}0.25}
\newcommand{\tfCorrUst}{0.20}
\newcommand{\tfCovidEq}{\ensuremath{-}20.2}
\newcommand{\tfCovidEqAbs}{20.2}
\newcommand{\tfCovidP}{11.9}
\newcommand{\tfFeeDrag}{2.8}
\newcommand{\tfFeeExact}{2.75}
\newcommand{\tfFeeInc}{0.75}
\newcommand{\tfFeeMgmt}{2.01}
\newcommand{\tfFirst}{January 2001}
\newcommand{\tfGainMgmt}{3.7}
\newcommand{\tfGfcEq}{\ensuremath{-}50}
\newcommand{\tfGfcEqAbs}{50}
\newcommand{\tfGfcP}{36}
\newcommand{\tfGocFirst}{January 2002}
\newcommand{\tfGocFirstPos}{February 2002}
\newcommand{\tfIncShareGain}{20}
\newcommand{\tfIncShareLeft}{37}
\newcommand{\tfLast}{August 2026}
\newcommand{\tfLeftMgmt}{1.99}
\newcommand{\tfLevMax}{3.1}
\newcommand{\tfLevMed}{2.3}
\newcommand{\tfMarkets}{7}
\newcommand{\tfMdd}{30}
\newcommand{\tfMddEq}{51}
\newcommand{\tfMean}{4.0}
\newcommand{\tfMeanNet}{1.2}
\newcommand{\tfMonths}{308}
\newcommand{\tfRfMean}{1.75}
\newcommand{\tfRhoAvg}{0.15}
\newcommand{\tfSR}{0.35}
\newcommand{\tfSRCad}{0.14}
\newcommand{\tfSREq}{0.43}
\newcommand{\tfSREur}{\ensuremath{-}0.01}
\newcommand{\tfSRGoc}{0.20}
\newcommand{\tfSRJpy}{0.37}
\newcommand{\tfSROil}{\ensuremath{-}0.04}
\newcommand{\tfSRUst}{0.21}
\newcommand{\tfSRavg}{0.18}
\newcommand{\tfSReq}{0.55}
\newcommand{\tfSRnet}{0.11}
\newcommand{\tfSRpred}{0.35}
\newcommand{\tfTurn}{7.2}
\newcommand{\tfTwentyTwoEq}{\ensuremath{-}20}
\newcommand{\tfTwentyTwoEqAbs}{20}
\newcommand{\tfTwentyTwoP}{27}
\newcommand{\tfTwentyTwoUst}{\ensuremath{-}16}
\newcommand{\tfTwentyTwoUstAbs}{16}
\newcommand{\tfVolCad}{8.2}
\newcommand{\tfVolEq}{15.6}
\newcommand{\tfVolEur}{9.1}
\newcommand{\tfVolGoc}{6.2}
\newcommand{\tfVolJpy}{9.5}
\newcommand{\tfVolNet}{11.4}
\newcommand{\tfVolOil}{39.1}
\newcommand{\tfVolP}{11.5}
\newcommand{\tfVolUst}{7.5}
\newcommand{\tfWorstTenEq}{\ensuremath{-}10.8}
\newcommand{\tfWorstTenP}{4.2}
\newcommand{\tfWorstTenUp}{7}
\newcommand{\tfWorstYear}{\ensuremath{-}15}
\newcommand{\tfWorstYearLab}{2016}
\newcommand{\tfYears}{25}
\newcommand{\tfYearsUp}{15}

\begin{refsection}
\setcounter{chapter}{8}
\chapter{Hedge Funds, Managed Futures and Digital Assets}
\label{ch:hedge}

An investment committee that adds hedge funds, a trend following programme or a sleeve of bitcoin to
a portfolio of stocks and bonds is buying something it cannot fully see. The manager's positions
change from month to month, the published track records come from databases that keep the winners
and lose the losers, and the newest assets have a history that is short, violent and chosen after the
fact. This chapter asks what each of these holdings is, what its record can and cannot show, and what
it did to a conventional portfolio. It ends with a test on public data: did a trend fund, after a
standard fee, improve a portfolio of 60\% US equities and 40\% US Treasury bonds over the twenty five
years from 2001, and did a small position in bitcoin improve it from 2015?

The chapter starts with the main families of hedge fund strategy and the kind of risk each one sells.
Trend following gets two sections, because it is a strategy whose rules can be stated explicitly and
tested on public prices without relying on a manager's reported returns: first what a trend rule can earn in theory, then a programme on seven markets
built from public data. A section on carry and roll yield adds the part of a futures return that the
programme leaves out. The middle of the chapter is about evidence. Fund databases overstate returns in
ways that can be computed, a regression splits a return into parts in a way that depends on the
factors chosen, and a regression clone of a fund shows how much of its return a few liquid positions
could have delivered. Then come fees: hurdles, high-water marks and the option value of an incentive
fee, and the second layer of a fund of funds, whose fees do not net across managers. The last part
turns to digital assets: the arithmetic of the protocol, volatility and correlation and how they have
changed, the selection that went into studying the survivors, and the size of position at which
bitcoin carries a given share of a portfolio's risk.

\begin{keybox}[How to read this chapter]
The numbers in the text are computed by thirteen scripts in \texttt{code/ch09/} (\texttt{trend},
\texttt{theory}, \texttt{carry}, \texttt{bias}, \texttt{survival}, \texttt{replicate},
\texttt{decompose}, \texttt{fees}, \texttt{fof}, \texttt{crypto}, \texttt{onchain},
\texttt{case\_portfolio} and \texttt{exercises}), and the tests in \texttt{code/ch09/test\_*.py} check
the library against simulation and numerical integration. Five short listings in
\texttt{code/ch09/snippets/} (\texttt{tsmom.py}, \texttt{permutation.py}, \texttt{rolldown.py},
\texttt{fees.R} and \texttt{barrier.m}) are printed in the text; a sixth, \texttt{issuance.m}, is
cited in \cref{sec9:crypto} but not printed; each prints the numbers quoted beside
it, and \texttt{code/ch09/test\_assets.py} checks them. An independent R script rebuilds the seven
market series from the daily files, runs the trend programme and its fee with its own loops, fits the
quarterly regression with \texttt{lm}, and recomputes the database formulas, the closure formulas
by numerical integration, the carry splits from the raw rate files, the attribution regression, the
fee formulas, Nakamoto's attacker table, the Fisher test, the fund of funds fees, the bitcoin statistics
and the risk shares. GNU Octave recomputes the normal formulas, the hurdle and closure formulas by
quadrature, the issuance sum, the risk share weight with \texttt{fzero} and the case study statistics
from the saved programme returns.

The trend programme uses monthly returns from \tfFirst{} to \tfLast, \tfMonths{} months. US equities are
the market excess return of the Kenneth R. French Data Library \parencite{frenchlib}, which also gives
the one month bill rate. The US 10 year yield, the price of WTI crude oil and the exchange rates of the
euro, the yen and the Canadian dollar against the US dollar are FRED series \parencite{fredch9}. The
Canadian 10 year benchmark yield and 90 day bill rate are from the Bank of Canada \parencite{bocvalet};
the carry section adds the 7 year benchmark yield and the one month bill yields of Canada and the US
used in Chapter~5. The programme's returns are built from these spot prices and yields, so they
approximate the returns of a futures programme and cannot be bought as an index.
Bitcoin and ether are the daily Coinbase prices published on FRED \parencite{coinbasefred}. The bitcoin
series starts on \btSeriesStart, and each month is sampled at its last available price, from
\btFirstPrice{} to \btLast, which gives monthly returns from \btFirst{} (ether from \ceFirst). The hedge funds of the database and replication sections
are simulated, and are labelled as such wherever they appear.

Random draws come from NumPy's default generator with base seed \texttt{9473} and a stream number
fixed in each script. The seed is an arbitrary label that makes the runs repeatable; it carries no
date or other meaning.

Notation. $r_t$ is a monthly excess return, the return over the one month bill, and $w_t$ a position
set at the end of month $t$ and held through month $t+1$. $\theta$ is an annual Sharpe ratio, $\sigma$
a volatility, $\varrho$ a correlation, and $\Phi$ and $\varphi$ are the standard normal distribution and
density. A fund's gain is its return over a year before incentive fees, and $x^+=\max(x,0)$.
\end{keybox}

\FloatBarrier
% =====================================================================================
\section{What a hedge fund sells}
\label{sec9:families}

A hedge fund is a pooled vehicle with few limits on what it holds: it can sell short, borrow, trade
derivatives and hold assets that are hard to value. Its fees, a management fee on assets and an
incentive fee on gains above a high-water mark, were the subject of Chapter~2. The name says less
than it seems to. Many hedge funds hedge only part of their market exposure, and some hedge none of
it. A more useful description of a strategy says which risk it takes on and in which states of the
world it loses money. \Cref{tab9:families} sorts the main families that way.

\begin{table}[t]
  \centering
  \small
  \caption{The main hedge fund strategy families, the source of their return and the shape of
  their payoff against the risk they carry. The descriptions are typical, and single funds differ.}
  \label{tab9:families}
  \setlength{\tabcolsep}{4pt}
  \begin{tabular}{@{}p{0.2\textwidth}p{0.4\textwidth}p{0.34\textwidth}@{}}
    \toprule
    Family & Main source of return & Payoff against its main risk\\
    \midrule
    Equity long/short & stock selection, plus a partial equity beta & roughly linear in the market\\
    Event driven & merger spreads, distressed debt & often like a short put on the deal or the market\\
    Relative value & convergence of close substitutes, carry & often like a short put on liquidity\\
    Global macro & views on rates, currencies and indices & varies with the manager\\
    Managed futures & trends in futures prices & often convex in sustained moves, like a long straddle\\
    Multi-strategy & a mix of the above & a mix of the above\\
    \bottomrule
  \end{tabular}
\end{table}

Three features recur across the table. First, much of the return is a payment for bearing a risk that
can be named and that public markets also price. An equity long/short fund with a beta of one half
earns half the equity premium before any skill. A merger arbitrage fund buys the target after a bid
is announced and earns the spread if the deal closes; \textcite{mitchell2001} show that its returns
behave like those of a short put on the market, with small steady gains and large losses in the months
when the market falls sharply and deals fail. Second, the payoffs are often nonlinear. Strategies that
earn a carry or a spread resemble the sale of insurance, and \textcite{agarwal2004} find that the
returns of many hedge fund indices load on the returns of out of the money index puts. Trend
followers are the exception that buys convexity, and \textcite{fung2001} show that their returns look
like those of lookback straddles. Third, many of the assets are illiquid, and a fund that marks them
from appraisals or stale quotes reports returns that are smoothed. Chapter~8 showed what smoothing
does to measured volatility and Sharpe ratios, and Chapter~7 that it hides beta; \textcite{asness2001}
find that adding lagged market returns to the regression raises the measured equity beta of hedge
fund indices substantially.

The rest of the chapter takes these points one at a time. Managed futures come first, because their
rules can be written down and run on public prices, so their record can be rebuilt rather than taken
from a database.

\FloatBarrier
% =====================================================================================
\section{What a trend rule earns}
\label{sec9:theory}

A \emph{time series momentum} rule, in the form studied by \textcite{moskowitz2012}, holds each market
long when its own return over the past year was positive and short when it was negative. The sign of
the position depends only on the market's own past, and the size is set so that each market carries a
similar risk. Before building such a programme it helps to know what the rule can earn and from where.
Three results answer that: one for the sign rule under normality, one for a market whose returns are
unpredictable, and an identity that shows why the payoff of a trend rule looks like that of an option.

\begin{proposition}[The sign rule]\label{prop9:sign}
Let $X$ and $Y$ be jointly normal with mean zero, standard deviations $\sigma_X$ and $\sigma_Y$, and
correlation $\varrho$. Then
\begin{equation}\label{eq9:sign}
  \E[\operatorname{sign}(X)\,Y] = \sqrt{2/\pi}\;\varrho\,\sigma_Y .
\end{equation}
For any $X$ and $Y$ with finite second moments,
\[
  \E[\operatorname{sign}(X)Y]=\E[\operatorname{sign}(X)]\,\E[Y]+\Cov(\operatorname{sign}(X),Y).
\]
\end{proposition}

\begin{proof}
Write $Y=\varrho\sigma_Y X/\sigma_X+\varepsilon$, where $\varepsilon$ is normal, independent of $X$ and has
mean zero. Then $\E[\operatorname{sign}(X)\varepsilon]=\E[\operatorname{sign}(X)]\E[\varepsilon]=0$ and
$\E[\operatorname{sign}(X)X]=\E|X|=\sigma_X\sqrt{2/\pi}$, so
$\E[\operatorname{sign}(X)Y]=\varrho\sigma_Y\sqrt{2/\pi}$. The second statement is the definition of the
covariance.
\end{proof}

Take $X$ to be the past year's return and $Y$ the next month's. A sign rule scaled to one unit of the
market's volatility then has an annual Sharpe ratio of about $\sqrt{12}\sqrt{2/\pi}\,\varrho$. The
approximation takes the rule's standard deviation to be $\sigma_Y$; the exact one is
$\sqrt{\sigma_Y^2-m^2}$, with $m$ the rule's mean, which is very close when $\varrho$ is small. A correlation of
0.05 between the past year and the next month is then worth a Sharpe ratio of about 0.14 per market. The
second statement of the proposition splits what the rule earns into two parts. The first,
$\E[\operatorname{sign}(X)]\E[Y]$, is a reward for being long on average in a market that drifts up; it
needs no predictability at all. The second, the covariance, is the timing part, and only it measures
whether the past says anything about the future.

\begin{proposition}[No trend, no profit]\label{prop9:mart}
Suppose that $\E[r_{t+1}\mid\mathcal F_t]=0$ for every $t$, where $\mathcal F_t$ is the information
available at the end of month $t$, and let $w_t$ be any position chosen with that information, with
$\E|w_tr_{t+1}|<\infty$. Then $\E[w_tr_{t+1}]=0$. If instead the returns are independent with a common
mean $\mu$, then $\E[w_tr_{t+1}]=\mu\,\E[w_t]$.
\end{proposition}

\begin{proof}
By the tower property, $\E[w_tr_{t+1}]=\E\bigl[w_t\,\E[r_{t+1}\mid\mathcal F_t]\bigr]=0$. When the returns
are independent with mean $\mu$ and $w_t$ depends only on returns up to $t$,
$\E[r_{t+1}\mid\mathcal F_t]=\mu$ and the same step gives $\mu\E[w_t]$.
\end{proof}

A trend programme that trades a market whose excess returns are a martingale difference earns
nothing before costs and loses its costs after them, however clever the rule. Whatever it earns
beyond the drift term must come from dependence between past and future returns. This gives a test
that needs no model: shuffle the months of the data and rerun the programme. A shuffle keeps each
month's returns across markets together but randomises their order in time, which breaks the link
between past and future returns that a trend rule could exploit. The shuffled runs still earn the
drift term, $\mu\,\E[w_t]$ in \cref{prop9:mart}, so what the test measures is timing.

\begin{proposition}[The straddle identity]\label{prop9:straddle}
Let $S_t=r_1+\dots+r_t$ with $S_0=0$, and hold the position $w_{t-1}=S_{t-1}$ through period $t$, so that
the position grows with the move since the start. Then
\begin{equation}\label{eq9:straddle}
  \sum_{t=1}^n S_{t-1}r_t = \tfrac12\Bigl(S_n^2-\sum_{t=1}^n r_t^2\Bigr).
\end{equation}
If the $r_t$ are independent and identically distributed with mean zero and variance $\sigma^2$, the expected profit is zero and the
probability of a loss tends to $2\Phi(1)-1\approx0.683$ as $n\to\infty$.
\end{proposition}

\begin{proof}
$S_t^2=S_{t-1}^2+2S_{t-1}r_t+r_t^2$; summing from $t=1$ to $n$ and using $S_0=0$ gives \cref{eq9:straddle}.
The expectation of the right side is $\tfrac12(n\sigma^2-n\sigma^2)=0$. A loss occurs when
$S_n^2/(n\sigma^2)<\sum r_t^2/(n\sigma^2)$. The right side tends to one by the law of large numbers and
$S_n/(\sigma\sqrt n)$ tends in law to a standard normal $Z$ by the central limit theorem, so the
probability tends to $\Pr(Z^2<1)=2\Phi(1)-1$.
\end{proof}

This identity appears as equation (7) of \textcite{dao2017}, who use it to read the profit of a trend
rule as long horizon variance less short horizon variance. The identity is exact for any path. Its right side is the payoff of a straddle on the total move,
$\tfrac12S_n^2$, paid for with the realised variance of the path, $\tfrac12\sum r_t^2$. The linear trend
rule therefore has a convex payoff in the total move, though it is not a purchased straddle and its
profit can be negative: it gains on large moves in either direction and loses a little, the realised
variance, on paths that go nowhere. On \stNpath{} simulated years of \stNday{} daily returns with mean zero, the two sides of
\cref{eq9:straddle} agree to within a relative error of \stMaxErr, rounding error only. The rule loses
money in \stLossShare\% of the years, with a simulation standard error of \stLossSe{} points. For normal
returns with mean zero the exact share follows from \cref{eq9:straddle}: since
$\sum r_t^2=S_n^2/n+\sum(r_t-\bar r)^2$, the rule loses exactly when $\lvert t\rvert<1$, where $t$ is
the usual $t$ statistic of the mean of the $n$ returns, so the probability is
$\Pr(\lvert t_{n-1}\rvert<1)$, \stLossFinite\% for $n=\stNday$, against the limit of
\stLossTheory\%; the simulated share is a draw about \stLossZ{} standard errors low on this seed. Its profits have a skewness
of \stSkew, against $2\sqrt2\approx2.8$ in the limit, where the profit is a multiple of $Z^2-1$. Panel (a) of \cref{fig9:straddle} shows the parabola.

\begin{figure}[!tb]
  \centering
  \includegraphics[width=\textwidth]{ch09/fig03_straddle}
  \caption{(a) Profit of the linear trend rule of \cref{prop9:straddle} against the total move over a
  simulated year of daily returns, with the parabola $\tfrac12(S^2-n\sigma^2)$. (b) Annual Sharpe ratios,
  after trading costs and before fees, of the trend programme of \cref{sec9:programme} on \pmN{} samples with the months shuffled, and the
  ratio on the data in order (vertical line).}
  \label{fig9:straddle}
\end{figure}


\paragraph{What the data say about the sign rule}
For each of the seven markets of the next section, the correlation between the past year's return and
the next month's return is small: \srRhoEq{} for US equities, \srRhoJpy{} for the yen and \srRhoOil{} for
oil, and \srRhoAvg{} on average. \Cref{eq9:sign}, which has no drift in it, would then predict a timing
gain of \srPredSRavg. The rule in fact earned a Sharpe ratio of \srRealSRavg{} per market on average. Of that,
the drift term of \cref{prop9:sign} accounts for \srDriftAvg{} and the timing term for \srTimingAvg; for
US equities the two are \srDriftEq{} and \srTimingEq. A timing gain with almost no linear correlation
means that the dependence is not of the joint normal kind \cref{eq9:sign} assumes. One possible reason,
which would need a test of its own, is that large moves are followed by moves of the same sign more
often than small ones, and a correlation, dominated by the many small months, hardly registers that.
\Cref{fig9:deciles} looks for that pattern and does not find it. Pooled over the seven markets, with
each series in units of its own standard deviation, the \dqObs{} market months have a correlation of
\dqRho. The bottom decile of past returns was followed by an average return of \dqBottom{} standard
deviations, a reversal, and the top decile by \dqTop; the months in between averaged \dqMid. The two
extreme deciles together took \dqTailShareAbs\% off the sign rule's pooled average, so what the rule
earned came from the middle of the distribution. The signals overlap from month to month and the
markets move together, so the decile means are not independent and the picture describes the sample
without testing anything.

\begin{figure}[!tb]
  \centering
  \includegraphics[width=0.74\textwidth]{ch09/fig17_deciles}
  \caption{Mean next-month return by decile of the past 12 month return, pooled over the seven markets,
  \tfFirst{} to \tfLast, each series divided by its own standard deviation, with the least-squares line
  through all \dqObs{} market months. The bars are placed at the mean signal of each decile.}
  \label{fig9:deciles}
\end{figure}


The shuffle test of \cref{prop9:mart} asks whether the programme as a whole earned more than chance.
The null is that the months are exchangeable. On \pmN{} shuffles, which keep each market's drift and
remove any dependence over time, the Sharpe ratio after trading costs and before fees has mean \pmMean{} and standard deviation \pmSd,
and its 95th percentile is \pmQ. Before trading costs the shuffled mean is \pmMeanGross, what the drift
earns, and the costs take most of that. The ratio on the data in order, \tfSR, is matched or exceeded by
\pmCount{} of the shuffles, a $p$-value of $(1+\pmCount)/(\pmN+1)=\pmP$ (\cref{fig9:straddle}, panel (b)). Under the exchangeability null this is significant at the 5\% level, but one quarter century of
seven markets would not persuade a sceptic of a stable timing advantage; that is why the longer
histories of \textcite{hurst2017}, which go back more than a century, carry more weight.

\FloatBarrier
% =====================================================================================
\section{A trend programme on seven markets}
\label{sec9:programme}

\paragraph{The markets}
The programme trades seven markets: US equities, the US and Canadian 10 year government bonds, WTI
crude oil, and the euro, yen and Canadian dollar against the US dollar. A professional programme
would trade futures on these and many more. Here each market's monthly excess return is built from
public month-end prices. US equities are the French market excess return. A bond return is that of a
par bond bought at the month-end yield $y_0$ with ten years to run and sold a month later at the new
yield $y_1$, with semiannual coupons:
\begin{equation}\label{eq9:bond}
  r^{\mathrm{bond}} = P\bigl(y_0,y_1,10-\tfrac1{12}\bigr)-1+\frac{y_0}{12},\qquad
  P(c,y,m)=\frac cy\bigl(1-(1+y/2)^{-2m}\bigr)+(1+y/2)^{-2m},
\end{equation}
less the one month US bill rate for the US bond and less one twelfth of the 90 day Canadian bill rate
for the Canadian bond, so that each is the return of a position financed in its own currency. The Canadian bond excess
return is measured in Canadian dollars; treating it as a US-dollar portfolio contribution assumes the
Canadian-dollar exposure is hedged, with forward-hedging costs and basis omitted from this proxy. The
currency returns are spot changes only, so the interest differential that a currency future would
earn or pay is left out, and the oil return is the change in the spot price, which leaves out the roll
yield of Chapter~4. Both omissions change the level of returns in those markets more than they change
the trend signals, but with them the programme is a proxy built from spot prices and par yields, and
its returns are those of no traded futures position. The bond formula is itself a yield-based
approximation: it prices a par bond with a fractional maturity, and the actual semiannual coupon dates,
accrued interest and settlement are not modelled. The Canadian bond's first signal is set at the end of \tfGocFirst, after a year of returns, so it
first holds a position in \tfGocFirstPos.

\paragraph{The rule}
At the end of each month $t$ the programme computes, for each market, the compounded excess return over
the past 12 months and the annualised standard deviation $\hat\sigma_{i,t}$ of up to 24 past monthly
returns (at least 12). It then holds through month $t+1$
\begin{equation}\label{eq9:weights}
  w_{i,t} = \operatorname{sign}\Bigl(\prod_{s=t-11}^{t}(1+r_{i,s})-1\Bigr)\,\frac{0.20}{N_t\,\hat\sigma_{i,t}},
\end{equation}
where $N_t$ is the number of markets with a signal, and pays 5 basis points on every unit of notional
traded. Each market therefore targets an annual volatility of 20\% divided by $N_t$. The split is made
before allowing for the correlations between markets, so it is not a volatility target for the
programme as a whole. The positions are
notional exposures financed at the bill rate, as futures are; the median sum of absolute positions is
\tfLevMed{} times capital and the largest \tfLevMax, and the turnover is \tfTurn{} times capital a year.

The listing below is the rule in pandas. The line that sets \texttt{held} is the one-month lag: a
position computed from returns up to the end of month $t$ earns the return of month $t+1$, and
dropping that shift is the easiest way to give a backtest information it could not have had. Run from
the repository root it prints Sharpe ratios of \lbSROne, \lbSRThree{} and \lbSRTwelve{} for lookbacks
of one, three and twelve months.
\Needspace*{32\baselineskip}
\lstinputlisting[language=Python]{code/ch09/snippets/tsmom.py}

The shuffle test of \cref{prop9:mart}, reported in \cref{sec9:theory}, is a loop around this
function. Each shuffle moves whole months, so each month's cross-section and each market's drift are
kept and only the order in time is lost; the \pmHeld{} months from \pmHeldFirst{} to \pmHeldLast, before the Canadian bond
series has a return, stay where they are, and only the \pmLive{} months with all seven returns are shuffled, so every shuffled
month keeps its whole cross-section. This is a permutation test under exchangeability of whole-month
return vectors; serial dependence in volatility or other time-varying features can invalidate its
nominal significance. On the chapter's random stream for the test it finds \pmCount{} of \pmN{} shuffles at least
as good as the data and $p=\pmP$. The 1 added to the numerator and the denominator counts the data in order as one of the
equally likely orderings, so the $p$-value cannot be zero.
\Needspace*{20\baselineskip}
\lstinputlisting[language=Python]{code/ch09/snippets/permutation.py}

\paragraph{Results}
From \tfFirst{} to \tfLast{} the programme earned an excess return of \tfMean\% a year with a volatility
of \tfVolP\%, a Sharpe ratio of \tfSR{} before fees, against \tfSReq{} for US equities over the same
months (\cref{fig9:trend}). Its correlation with equities was \tfCorrEq{} and its beta \tfBetaEq. Its
worst drawdown was \tfMdd\%, against \tfMddEq\% for equities, and it rose in \tfYearsUp{} of the
\tfYears{} full calendar years, with the best year \tfBestYearLab{} (\tfBestYear\%) and the worst
\tfWorstYearLab{} (\tfWorstYear\%).

\begin{figure}[!tb]
  \centering
  \includegraphics[width=\textwidth]{ch09/fig01_trend}
  \caption{(a) Growth of one unit invested in the excess return of the trend programme, before and
  after a fee of 2\% of assets and 20\% of gains above a high-water mark, and of US equities, \tfFirst{}
  to \tfLast, log scale. (b) Drawdowns of the programme before fees and of US equities.}
  \label{fig9:trend}
\end{figure}


The record in the worst months for equities is the reason investors hold such programmes. From
November 2007 to February 2009, when US equities lost \tfGfcEqAbs\% in total return, the programme
returned \tfGfcP\% including the bill rate. In 2022, when equities lost \tfTwentyTwoEqAbs\% and the US 10
year bond \tfTwentyTwoUstAbs\%, it returned \tfTwentyTwoP\%, and in February and March 2020 it gained \tfCovidP\% while equities lost
\tfCovidEqAbs\%. These are total returns, after the programme's trading costs and before any fee; the
monthly figures that follow are excess returns. In the ten worst months for equities, which
averaged \tfWorstTenEq\%, it averaged \tfWorstTenP\% and rose in \tfWorstTenUp{} of them. A
twelve month signal turns slowly, and a crash that follows a long rise finds the programme long; the
protection comes in falls that last long enough for the signals to turn.
\Cref{fig9:positions} shows the positions month by month. The programme was short equities in
\psGfcShortEq{} of the \psGfcMonths{} months from November 2007 to February 2009, from
\psGfcFirstShort{} on. It went into February and March 2020 long equities, with positions of
\psEqFebTwenty{} and \psEqMarTwenty{} times capital, and held a short position of \psEqAprTwenty{} times
capital only in April. Through 2022 it was short both bonds in all \psTwoTwoShortUst{} months.

\begin{figure}[!tb]
  \centering
  \includegraphics[width=0.96\textwidth]{ch09/fig15_positions}
  \caption{Positions held by the programme in each month, \tfFirst{} to \tfLast, in units of capital, long
  in blue and short in red. The boxes mark November 2007 to February 2009, February and March 2020, and
  2022.}
  \label{fig9:positions}
\end{figure}


\paragraph{The smile}
\Cref{prop9:straddle} suggests that a trend programme's return should be a convex function of the
move in the markets it trades. \Cref{fig9:smile} plots the programme's return in each of \slQuarters{}
calendar quarters against the equity excess return in the same quarter. A regression of the first on
the second and its square gives a slope of \slLin{} ($t=\slLinT$) and a coefficient on the square of
\slQuad{} ($t=\slQuadT$), with returns in decimals, and an intercept of \slIcpt\% a quarter. The
positive coefficient on the square says that the relation was convex in this sample; it is too
imprecise to promise the same protection in another. In the fifth of quarters with the
lowest equity returns the programme averaged \slLowP\%, in the middle three fifths \slMidP\% and in the
best fifth \slHighP\%. The shape is that of \textcite{fung2001}, with a weaker right arm: the programme
gained most when equities fell hard, and little when they rose hard.

\begin{figure}[!tb]
  \centering
  \includegraphics[width=0.72\textwidth]{ch09/fig02_smile}
  \caption{Quarterly return of the trend programme before fees against the US equity excess return in
  the same quarter, \slQuarters{} quarters, with the fitted quadratic.}
  \label{fig9:smile}
\end{figure}


\paragraph{The lookback and the markets}
\Cref{fig9:lookback} repeats the programme with lookbacks from one to 24 months. The Sharpe ratio is
\lbSROne{} at one month, \lbSRThree{} at three, \lbSRNine{} at nine, \lbSRTwelve{} at twelve and
\lbSRTwentyFour{} at 24. The standard error of an annual Sharpe ratio estimated from \tfMonths{} months
is about \lbSE. The variants are highly correlated, so their differences are estimated more precisely
than that, but not precisely enough to make nine months a better choice than twelve, and picking the
best lookback after the fact would overstate what a manager could have expected. Market
by market, the sleeves earned Sharpe ratios from \tfSROil{} for oil and \tfSREur{} for the euro to
\tfSRJpy{} for the yen and \tfSREq{} for US equities.

\begin{figure}[!tb]
  \centering
  \includegraphics[width=\textwidth]{ch09/fig04_lookback}
  \caption{(a) Sharpe ratio of the programme before fees for lookbacks of 1 to 24 months. (b) Sharpe
  ratio of each market's sleeve with the 12 month lookback, over the months the sleeve was active.}
  \label{fig9:lookback}
\end{figure}


The programme's ratio is about twice the average of its sleeves, and a short calculation says why.

\begin{proposition}[Diversification across markets]\label{prop9:diversify}
Let $n$ strategies have the same expected excess return $m$, the same volatility $s$, so the same
Sharpe ratio $\theta=m/s$, and pairwise correlation $\varrho$. The equally weighted combination has
Sharpe ratio
\begin{equation}\label{eq9:diversify}
  \theta_n = \theta\sqrt{\frac{n}{1+(n-1)\varrho}},
\end{equation}
which rises with $n$ to the limit $\theta/\sqrt\varrho$ when $\varrho>0$.
\end{proposition}

\begin{proof}
The combination has mean $m$ and variance $n^{-2}\bigl(ns^2+n(n-1)\varrho s^2\bigr)=s^2\bigl(1+(n-1)\varrho\bigr)/n$.
Dividing gives \cref{eq9:diversify}; the limit follows by dividing numerator and denominator by $n$.
\end{proof}

The seven sleeves have an average Sharpe ratio of \tfSRavg, each over its own active months, and an
average pairwise correlation of \tfRhoAvg{} over the \scMonths{} months from \tfGocFirstPos{} to
\tfLast{} when all seven were active. The \tfSRavg{} is above the \srRealSRavg{} of the one-unit sign
rule per market in \cref{sec9:theory} because each sleeve sizes its position by the inverse of its
trailing volatility and pays trading costs. \Cref{eq9:diversify} then gives \tfSRpred, against
the programme's \tfSR. The agreement is closer than such a rough formula deserves, since the sleeves
have neither equal ratios nor equal risks, but the lesson is general: diversification across
imperfectly correlated markets lifts a trend programme's ratio well above that of its average market,
and no single market accounts for its quality. It
is also why the many markets of a commercial programme matter, and why a seven market version is a
modest stand-in.
\Cref{fig9:sleevecorr} shows how far the sleeves are from one common correlation. Over those
\scMonths{} months the pairwise correlations run from \scMin{} for \scMinPair{} to \scMax{} for
\scMaxPair, with a standard deviation of \scSdPairs{} across the \scPairs{} pairs. The two bond
sleeves behave almost as one market, and \scNeg{} pairs are negatively correlated.

\begin{figure}[!tb]
  \centering
  \includegraphics[width=0.6\textwidth]{ch09/fig16_sleevecorr}
  \caption{Correlations of the monthly returns of the seven sleeves over the \scMonths{} months, \tfGocFirstPos{}
  to \tfLast, in which all seven held positions. The colour scale runs from $-0.7$ to $0.7$; the diagonal is left blank.}
  \label{fig9:sleevecorr}
\end{figure}


\paragraph{After fees}
A fee of 2\% of assets a year, charged monthly, and 20\% of gains above a high-water mark,
crystallised each December and at the end of the record, takes \tfFeeDrag\% a year from the programme. The net excess return falls
to \tfMeanNet\% and the Sharpe ratio from \tfSR{} to \tfSRnet{} (\cref{fig9:trend}). The management fee
is the larger part, \tfFeeMgmt{} of the \tfFeeExact{} points before rounding. A programme run at a
volatility of \tfVolP\% that earns \tfMean\% a year over bills before fees loses half of that to a fee
of 2\% of assets whatever happens. The incentive fee takes the other \tfFeeInc{} points. That is close
to \tfIncShareGain\% of the fund's gain after the management fee, \tfGainMgmt\% a year with the bill rate
of \tfRfMean\% included, but \tfIncShareLeft\% of the \tfLeftMgmt\% it earns over bills (both shares before
rounding), because it is charged on the whole gain above the high-water mark, bill rate included, in good years only, with
nothing returned in bad ones. The net volatility is \tfVolNet\%. Chapter~2 showed how much such a fee stream is worth to the
manager. For the investor the lesson is that a fixed fee on assets is heavy for a programme run at a
moderate volatility, since the fee does not shrink with the risk taken.
The R listing below applies the fee to the gross programme: the management fee first, every month,
then at each December the incentive fee on the whole gain above the mark, bill rate included. It
prints the split of \tfFeeMgmt{} and \tfFeeInc{} points and the Sharpe ratios of \tfSR{} and \tfSRnet,
and checks \cref{eq9:hwm} below on the whole record: the incentive fees paid equal $f/(1-f)$ times the
rise in the high-water mark.
\Needspace*{26\baselineskip}
\lstinputlisting[language=R]{code/ch09/snippets/fees.R}

\FloatBarrier
% =====================================================================================
\section{Carry and roll yield}
\label{sec9:carry}

A managed futures programme holds futures and forwards, and the return of a futures position is more
than the change in the price of the underlying. The difference is the \emph{carry}: the return the
position would earn if the spot price did not move, with the futures price converging to it. \textcite{koijen2018} define carry in this
way for every asset class and show that it predicts returns in many of them. For commodities the same
quantity is usually called the roll yield, and \textcite{erb2006} argue that differences in roll
returns account for much of the difference in long-run returns across commodity futures. The trend
programme above leaves out the carry of currencies and the roll of oil by working with spot prices,
and its bond returns include the coupon carry but not the roll-down. This
section shows what that leaves out, first in general and then for a currency and a bond.

\begin{proposition}[Futures return, spot return and carry]\label{prop9:carry}
Let $S_t$ be the spot price and $F_t$ the price at $t$ of a futures contract that expires at $t+1$, so
that $F_{t+1}=S_{t+1}$. Define the carry $C_t=(S_t-F_t)/F_t$ and the spot return $r^S=(S_{t+1}-S_t)/S_t$.
The return of a long futures position per unit of notional, which leaves out the interest on its
collateral and so is an excess return, is
\begin{equation}\label{eq9:carry}
  r^F=\frac{F_{t+1}-F_t}{F_t}=C_t+(1+C_t)\,r^S .
\end{equation}
If the futures price is set by the cost of carry, $F_t(\tau)=S_t\,e^{(r-y)\tau}$ for a contract with time
$\tau$ to expiry, with $r$ the financing rate and $y$ the yield of holding the underlying, then
$C_t=e^{(y-r)\Delta}-1$ for a contract with $\tau=\Delta$. A contract bought with time $\tau$ to expiry and
held for a period $\Delta$, while $r$ and $y$ stay constant, earns
\begin{equation}\label{eq9:roll}
  \log\frac{F_{t+\Delta}(\tau-\Delta)}{F_t(\tau)}=\log\frac{S_{t+\Delta}}{S_t}-(r-y)\Delta ,
\end{equation}
so a position rolled at the end of each period earns the log spot return plus $y-r$ per unit of time.
The simple return in \cref{eq9:carry} has the cross term $C_tr^S$; the additive split into spot return
and carry is exact for log returns under the cost of carry.
\end{proposition}

\begin{proof}
$(S_{t+1}-F_t)/F_t=(S_t-F_t)/F_t+(S_{t+1}-S_t)/F_t$, and $S_t/F_t=1+C_t$, which gives \cref{eq9:carry}. Under the
cost of carry $S_t/F_t(\Delta)=e^{(y-r)\Delta}$. For \cref{eq9:roll}, the contract bought at $t$ with time
$\tau$ to run has time $\tau-\Delta$ to run at $t+\Delta$, and
$\log F_{t+\Delta}(\tau-\Delta)-\log F_t(\tau)=\log S_{t+\Delta}+(r-y)(\tau-\Delta)-\log S_t-(r-y)\tau$.
\end{proof}

The yield $y$ is what holding the asset pays. For a currency it is the foreign interest rate, so the
carry of a long currency forward is the foreign rate less the domestic one; covered interest parity,
which sets the forward price, has held only approximately since 2008 \parencite{du2018}. For a bond it
is the coupon, and a bond also rolls down an upward sloping curve. For a commodity it is the
convenience of holding the physical good less the cost of storing it. A commodity curve that rises
with maturity, in contango, has negative carry: with spot at \oilS{} dollars and the one month future at
\oilFOne, \cref{eq9:roll} says a rolled long position earns, beyond the spot return, $12\log(\oilS/\oilFOne)$, about
\oilRoll\% a year in log terms, if the curve keeps its shape and is priced by the cost of carry. Positive carry is not a free
return. It is what the holder is paid, and the spot price may move against the holder by more.

\paragraph{The Canadian dollar}
Take an investor who holds one month Canadian Treasury bills funded by borrowing at the US one month
bill rate, the cash version of a long position in a Canadian dollar forward. Let $S$ be the price of a
US dollar in Canadian dollars and $i^{c}$ and $i^{u}$ the two bill yields, quoted as annual rates and
accrued at one twelfth a month, an approximation whose accuracy depends on how each yield is quoted.
The excess return in US dollars of the funded position over a month then splits exactly as
\begin{equation}\label{eq9:fxcarry}
  \Bigl(1+\tfrac{i^c}{12}\Bigr)\frac{S_t}{S_{t+1}}-\Bigl(1+\tfrac{i^u}{12}\Bigr)
  =\underbrace{\frac{S_t}{S_{t+1}}-1}_{\text{spot}}
  +\underbrace{\frac{i^c-i^u}{12}}_{\text{carry}}
  +\underbrace{\frac{i^c}{12}\Bigl(\frac{S_t}{S_{t+1}}-1\Bigr)}_{\text{cross}} ,
\end{equation}
which the code checks to within \cadIdGap{} a month. From \cadFirst{} to \cadLast, \cadMonths{} months, the
position earned \cadEx\% a year with a volatility of \cadVol\%: the spot part gave \cadSpot\%, the carry
\cadCarry\% and the cross term \cadCross\%. The carry was positive in \cadCarryPos\% of months and ranged
from \cadCarryMin\% to \cadCarryMax\% a year, small beside an exchange rate that moves by several percent
in a month (\cref{fig9:carry}, panel (a)). A regression of the monthly excess return on the carry
known at the start of the month has a slope of \cadSlope{} with a standard error of \cadSlopeSe. Uncovered
interest parity predicts a slope of zero, and a random walk in the exchange rate predicts one; one
currency pair over two decades cannot tell them apart. The evidence for currency carry comes from many
pairs at once \parencite{koijen2018}.

\begin{figure}[!tb]
  \centering
  \includegraphics[width=\textwidth]{ch09/fig13_carry}
  \caption{(a) Canadian bills funded at the US bill rate, \cadFirst{} to \cadLast: cumulative sums of the
  monthly excess return in US dollars and of its spot and carry parts in \cref{eq9:fxcarry}. (b) The
  Canada 10 year par bond over the 90 day bill, \gcFirst{} to \gcLast: cumulative sums of the realised
  excess return, of the carry, and of carry and roll-down together, all monthly.}
  \label{fig9:carry}
\end{figure}


\paragraph{The Canada 10 year bond}
For a par bond with yield $y_{10}$ financed at the bill rate $b$, the carry over a month is
$(y_{10}-b)/12$, and the roll-down is the gain from repricing the bond a month later at the yield of a 9
year 11 month bond on the same curve. The code interpolates that yield between the 7 and 10 year
benchmark yields of the Bank of Canada and prices the bond with the par bond formula. The expected
return on an unchanged curve, carry and roll-down together, is known at the start of each month. The
realised return prices the same bond at the next month's interpolated yield, so the difference between
the two is the effect of the change in the interpolated curve between 7 and 10 years, which the
quoted 10 year yield alone does not capture. From \gcFirst{} to \gcLast{} the carry averaged
\gcCarry\% a year and the roll-down \gcRoll\%, for an expected \gcExp\%; the realised excess return was
\gcReal\%. The remaining \gcRest{} percentage points is the realised-minus-expected return under the
interpolated-curve model; it reflects changes in the curve, over a period when the 10 year yield fell
from \gcYStart\% to \gcYEnd\%, and is not separately identified as the effect of that yield alone. The expected
return was positive in \gcExpPos\% of months. Its volatility was \gcVolExp\% a year against \gcVolReal\% for
the realised return, so in any one month the expected return is lost in the noise. Over twenty five
years carry alone gave \gcCarryShare\% of the bond's excess return, and carry and roll-down together
\gcExpShare\% (\cref{fig9:carry}, panel (b)). The slope of the realised
on the expected monthly return is \gcSlope, with a standard error of \gcSlopeSe.
The listing below does the calculation from the Bank of Canada file and prints the means of the
carry, \gcCarry\%, the roll-down, \gcRoll\%, the realised excess return, \gcReal\%, and the curve term,
\gcRest\%, in percent a year. The curve term is realised minus expected by construction, which is the
sense in which it measures the change in the interpolated curve.
\Needspace*{24\baselineskip}
\lstinputlisting[language=Python]{code/ch09/snippets/rolldown.py}

\FloatBarrier
% =====================================================================================
\section{What fund databases leave out}
\label{sec9:biases}

Most of what is known about hedge fund returns comes from commercial databases to which funds report
voluntarily. A fund starts reporting when it chooses, usually after a good start, and stops when it
chooses, often after losses or before closing. Three biases follow. \emph{Survivorship bias} arises when
the funds that stopped reporting are dropped from the sample. \emph{Backfill} or \emph{instant history bias}
arises when a fund that joins the database brings its earlier, and selected, history with it.
\emph{Selection bias} arises because the funds that report at all differ from those that do not.
\textcite{brown1992} showed for mutual funds how survivorship alone creates apparent persistence
in performance, \textcite{fung2000} and \textcite{liang2000} measured the first two in hedge fund
databases, and \textcite{aiken2013}, using returns of funds that do not report, found that the
funds missing from commercial databases did worse than those in them. \textcite{grecu2007} asked why funds
stop reporting and found that returns are markedly worse at the end of a fund's reporting life, and that
larger funds and better performing funds are less likely to fail, which points to failure rather than
success as the usual reason. \textcite{agarwal2013} used the equity holdings that funds
must disclose to compare those that report to databases with those that do not. Each bias can be computed under simple assumptions.

\begin{proposition}[The mean of the survivors]\label{prop9:survivor}
Let a fund's $T$ monthly returns be independent $\N(\mu,\sigma^2)$, and let it survive when its average
return $\bar R$ exceeds $c$. With $s=\sigma/\sqrt T$, $a=(c-\mu)/s$ and $\lambda(a)=\varphi(a)/(1-\Phi(a))$,
\begin{equation}\label{eq9:survivor}
  \E[\bar R\mid\bar R>c] = \mu + s\,\lambda(a),
\end{equation}
and the fraction that survives is $1-\Phi(a)$.
\end{proposition}

\begin{proof}
$\bar R$ is $\N(\mu,s^2)$. Write $\bar R=\mu+sZ$; then $\E[Z\mid Z>a]=\int_a^\infty z\varphi(z)\dd z/(1-\Phi(a))
=\varphi(a)/(1-\Phi(a))$, because $z\varphi(z)=-\varphi'(z)$.
\end{proof}

The bias $s\lambda(a)$ grows with the volatility of the funds and shrinks with the length of the
record, and it is positive for every cut-off. With $\mu=\svMu\%$ and $\sigma=\svSig\%$ a month and a
three year record, a rule that keeps funds with a positive average lets \svShareZero\% survive, and
their average exceeds the truth by \svBiasZeroTh\% a year. A simulation of 40,000 funds gives
\svBiasZero\%, and across cut-offs from minus 1\% to 1.2\% a month the formula and the simulation differ
by at most \svMaxGap{} basis points a month (\cref{fig9:survivor}). The bias is largest when the cut-off
removes most funds, which is the case for a database of funds that survived a long and difficult period.

\begin{figure}[!tb]
  \centering
  \includegraphics[width=0.72\textwidth]{ch09/fig05_survivor}
  \caption{Survivorship bias in the average return of funds with $T=\svT$ months of independent normal
  returns, $\mu=\svMu\%$ and $\sigma=\svSig\%$ a month, against the survival cut-off on the monthly mean:
  \cref{eq9:survivor} (line) and a simulation (points). The dashed line, on the right scale, is the share
  of funds that survive.}
  \label{fig9:survivor}
\end{figure}


\begin{proposition}[Backfill]\label{prop9:backfill}
Let a fund's returns be independent $\N(\mu,\sigma^2)$, and let it join the database after an
incubation of $K$ months only if its average over those months exceeds $c$, bringing that history with
it, and report $M$ further months. Among the funds that join, the average of the $M$ later months has
mean $\mu$, and the average of the whole record of $K+M$ months has mean
\begin{equation}\label{eq9:backfill}
  \mu + \frac{K}{K+M}\,\frac{\sigma}{\sqrt K}\,\lambda\Bigl(\frac{(c-\mu)\sqrt K}{\sigma}\Bigr).
\end{equation}
\end{proposition}

\begin{proof}
The later returns are independent of the incubation returns, so conditioning on the incubation leaves
their mean at $\mu$. The incubation average has the conditional mean of \cref{prop9:survivor} with
$T=K$, and the record average is the weighted mean of the two averages with weights $K/(K+M)$ and
$M/(K+M)$.
\end{proof}

With the same fund parameters, a two year incubation, entry when the incubation's compounded return
is positive and five years of later reporting, \bfEnterShare\% of funds join. Their incubation months
average \bfIncMean\% a year in simulation against \bfIncTheory\% from \cref{eq9:survivor}, and their
later months \bfAfterMean\%, the true mean. The formula selects on a positive arithmetic mean and the
simulation on a positive compounded return, so it is an approximation here. The full record is biased by
\bfRecordBias\% a year, against \bfRecordTheory\% from \cref{eq9:backfill} with the same approximation. The cure is simple when the
entry date is known: drop the months before it.

\paragraph{A simulated database}
The two biases act together in a database. The simulation behind \cref{fig9:backfill} starts
\dbFunds{} funds at random dates. Each has its own mean, drawn around 0.6\% a month with a standard
deviation of 0.2\%, and a monthly volatility of 3\%. A fund joins after a two year incubation with a
positive compounded return, and stops reporting at the first of three events: a trailing twelve month
return below minus 15\%, a random exit with a probability of 0.4\% a month, or ten years of reporting. Over a window of
\dbMonths{} months the database ends with \dbNdb{} funds that ever reported, of which \dbNlive{} still
report at the end; \dbDeadShare\% of those that reported stopped during the window. The true average
return of all fund months in the window is \dbTruth\% a year. The funds alive at the end average
\dbLive\%, a survivorship bias of \dbSurvBias\%. All reporting funds, dead ones included but with their
backfilled histories, average \dbBack\%; dropping the months before entry gives \dbClean\%, so the
backfill bias is \dbBackBias\% and \dbCleanBias\% remains, the differences computed before rounding. That remainder is the selection of which funds
enter, which no treatment of the database can undo.

\begin{figure}[!tb]
  \centering
  \includegraphics[width=0.72\textwidth]{ch09/fig06_backfill}
  \caption{Average return of the simulated funds by month relative to their entry into the database.
  The bars before entry are backfilled history; the dashed line is the true average of all fund
  months.}
  \label{fig9:backfill}
\end{figure}


In this simulation the backfilled months average \dbPreEntry\% a year and the months after entry,
averaged month by month over the \dbEventLast{} months that follow it as in \cref{fig9:backfill},
\dbPostEntry\%, against the \dbClean\% of all post-entry months pooled.
\Cref{fig9:backfill} shows the step at the entry date. The estimates for real
databases depend on the database and the period. \textcite{fung2000} put survivorship bias at a few
percent a year and instant history bias at somewhat over one percent, and \textcite{ackermann1999}
found smaller figures in a different sample. The simulation is not a calibration to any of them. It
shows that biases of this size arise from simple and plausible reporting rules, and that they can be
large beside the alpha a fund is supposed to deliver.

\paragraph{Funds that close after a deep loss}
Funds also leave because they fail. A simple rule is that a fund closes the first time it has lost a
fraction $\ell$ of its starting value, which makes survival depend on the whole path of returns. In
continuous time the survival probability and the survivors' mean have closed forms.

\begin{proposition}[Closure at a loss limit]\label{prop9:closure}
Let the log value of a fund be $X_t=\mu t+sW_t$, with $W$ a standard Brownian motion, and let the fund
close at the first time $X$ reaches $-b$, where $b=-\log(1-\ell)>0$. Write $m=\mu T$, $v=s\sqrt T$ and
$k=e^{-2\mu b/s^2}$. The probability that the fund is still open at $T$ is
\begin{equation}\label{eq9:closure}
  P_T=\Phi\Bigl(\frac{m+b}{v}\Bigr)-k\,\Phi\Bigl(\frac{m-b}{v}\Bigr),
\end{equation}
and the mean log return of the funds still open at $T$ is
\begin{equation}\label{eq9:closuremean}
  \E[X_T\mid \text{open at }T]=\frac1{P_T}\Bigl[m\,\Phi\Bigl(\tfrac{m+b}{v}\Bigr)+v\,\varphi\Bigl(\tfrac{m+b}{v}\Bigr)
  -k\Bigl((m-2b)\,\Phi\Bigl(\tfrac{m-b}{v}\Bigr)+v\,\varphi\Bigl(\tfrac{m-b}{v}\Bigr)\Bigr)\Bigr].
\end{equation}
\end{proposition}

\begin{proof}
By the reflection principle for Brownian motion with drift \parencite{shreve2004}, the paths that stay
above $-b$ up to $T$ and end near $x>-b$ have density
$v^{-1}\bigl[\varphi\bigl((x-m)/v\bigr)-k\,\varphi\bigl((x+2b-m)/v\bigr)\bigr]$. Integrating over $x>-b$ gives
\cref{eq9:closure}. For the mean, $\int_{-b}^\infty x\,v^{-1}\varphi((x-m)/v)\dd x=m\Phi((m+b)/v)+v\varphi((m+b)/v)$, and
the substitution $u=x+2b$ turns the second integral into $\int_b^\infty(u-2b)v^{-1}\varphi((u-m)/v)\dd u
=(m-2b)\Phi((m-b)/v)+v\varphi((m-b)/v)$, using $\varphi$ even. Dividing by $P_T$ gives \cref{eq9:closuremean}.
\end{proof}

\begin{figure}[!tb]
  \centering
  \includegraphics[width=\textwidth]{ch09/fig12_closure}
  \caption{Funds with no skill, a volatility of \clSig\% and closure after a loss of \clLoss\%. (a) The share
  still open: \cref{eq9:closure} with continuous monitoring, the same formula with the barrier moved for
  monthly monitoring, and a monthly simulation of \clFunds{} funds. (b) The survivors' mean log return
  per year, from \cref{eq9:closuremean} and the simulation; the true mean is zero.}
  \label{fig9:closure}
\end{figure}

A database checks a fund once a month, and a fund that dips below the limit between month ends is not
caught. \textcite{broadie1997} show, for barrier options, that monitoring at intervals $\Delta t$ is close to
continuous monitoring of a barrier moved away by $0.5826\,s\sqrt{\Delta t}$, a correction that becomes
exact as $\Delta t$ shrinks and is only approximate for monthly checks; the same correction applies here. Take funds with no skill,
$\mu=0$, a volatility of \clSig\% and a loss limit of \clLoss\%, so $b=\clB$, and simulate \clFunds{} of them
month by month for \clYears{} years (\cref{fig9:closure}). After five years \clSurvSimFive\% are open, against
\clSurvFive\% from \cref{eq9:closure} and \clSurvBgkFive\% with the correction; after ten years the figures are
\clSurvSimTen\%, \clSurvTen\% and \clSurvBgkTen\%, with a simulation standard error of \clSurvSe{} points. The
survivors' mean log return, which is zero for every fund before selection, is \clMeanSimFive\% a year
after five years and \clMeanSimTen\% after ten in the simulation, \clMeanBgkFive\% and \clMeanBgkTen\% with the
corrected formula. The bias grows with the horizon, because more of the losers have been removed, and
it comes with no sign in the survivors' own records that anything was removed.
The Octave listing evaluates \cref{eq9:closure}, checks it by integrating the density of the
survivors, the normal density less its image about the barrier, and then moves the barrier, not the
volatility, for monthly checks. It prints \clSurvFive\% and \clSurvBgkFive\% at five years and
\clSurvTen\% and \clSurvBgkTen\% at ten.
\Needspace*{21\baselineskip}
\lstinputlisting[language=Octave]{code/ch09/snippets/barrier.m}

\FloatBarrier
% =====================================================================================
\section{Where a fund's return comes from}
\label{sec9:decomp}

Before asking whether a fund can be copied, it helps to split its return into parts. There are two
standard splits. The first adds up the profit of each position and needs the fund's holdings. The
second regresses the fund's return on the returns of a set of factors and needs only the fund's
record; the intercept is called alpha and the rest is attributed to the factors. Least squares makes
the second split exact in sample, whatever factors are chosen.

\begin{proposition}[Return attribution by regression]\label{prop9:attrib}
Let $y_1,\dots,y_T$ be a fund's returns and $f_t\in\mathbb R^k$ factor returns, and fit
$y_t=\hat\alpha+\hat\beta^\top f_t+\hat e_t$ by least squares with an intercept. Then, with bars for sample
means and $\widehat{\Var}$ for sample variances,
\begin{equation}\label{eq9:attrib}
  \bar y=\hat\alpha+\hat\beta^\top\bar f,\qquad
  \widehat{\Var}(y)=\widehat{\Var}(\hat\beta^\top f)+\widehat{\Var}(\hat e).
\end{equation}
\end{proposition}

\begin{proof}
The normal equations say the residuals are orthogonal to the constant and to each factor:
$\sum_t\hat e_t=0$ and $\sum_t\hat e_tf_t=0$. Averaging the fitted equation and using the first gives the
mean identity. The second makes the residuals uncorrelated in sample with $\hat\beta^\top f_t$, so the
sample variance of the sum is the sum of the sample variances.
\end{proof}

The identity holds for any choice of factors, and that is its weakness: the split of the mean into
alpha and factor parts depends entirely on which factors are used. \Cref{tab9:decomp} applies both splits
to the trend programme of \cref{sec9:programme} over its \dcMonths{} months. The sleeve split is exact by
construction and shows where the \dcMean\% a year came from: most of it from US equities, the yen and the
two bond markets, with oil and the euro close to zero. The first regression uses the seven market
returns of the same month, over the \dcRegMonths{} months in which all seven have a return. Its $R^2$ is \dcRtwoA, its alpha \dcAlphaA\% a year with a $t$ statistic of \dcTA,
and the code confirms the mean identity of \cref{eq9:attrib} to within \dcIdGap. The $t$ statistic uses
ordinary least squares errors, so with serial dependence and changing volatility it is descriptive. The alpha is larger than
the mean return because the regression betas on US equities and the Canadian bond, two markets that
rose over the period, are negative.

\textcite{fung2001} explained trend followers' returns with lookback straddles, options whose payoff
depends on the range of prices within the month. A crude stand-in, the absolute return of each market
in the same month, adds seven factors that a straddle on each market would roughly pay. The $R^2$ rises
only to \dcRtwoB; of the seven coefficients, \dcAbsPos{} are positive and none has a $t$ statistic above two.
The alpha becomes \dcAlphaB\% a year with a $t$ statistic of \dcTB, and the part of the mean attributed
to the factors rises from \dcLoadPartA\% to \dcLoadPart\% a year, a shift of \dcLoadShift{} points, on almost
no change in fit. \Cref{prop9:straddle} is exact only for the linear rule; the programme's sign rule
shares its convexity in moves that build over the twelve month lookback without being the same payoff,
and a factor built from one month's move misses most of such a move. Both regressions illustrate \cref{eq9:attrib}: they fit equally badly and
attribute the same mean in opposite ways.

\begin{table}[t]
  \centering
  \caption{The trend programme, gross of fees, \tfFirst{} to \tfLast. Sleeve: the mean excess return of each
  market's positions, in percent a year, which add up to the programme's \dcMean\% before rounding. Market beta: the
  coefficient on the market's monthly return in the regression on the seven markets, over the
  \dcRegMonths{} complete months. Absolute move:
  the coefficient on the absolute monthly return of the market, and its $t$ statistic, when those seven
  are added.}
  \label{tab9:decomp}
  \small
  \begin{tabular}{@{}lrrrr@{}}
    \toprule
    Market & Sleeve, \% a year & Market beta & Absolute move & $t$ \\
    \midrule
    US equities & \dcSlEq & \dcBetaEq & \dcAbsEq & \dcAbsTEq \\
    US 10 year bond & \dcSlUst & \dcBetaUst & \dcAbsUst & \dcAbsTUst \\
    Canada 10 year bond & \dcSlGoc & \dcBetaGoc & \dcAbsGoc & \dcAbsTGoc \\
    WTI crude oil & \dcSlOil & \dcBetaOil & \dcAbsOil & \dcAbsTOil \\
    Euro & \dcSlEur & \dcBetaEur & \dcAbsEur & \dcAbsTEur \\
    Yen & \dcSlJpy & \dcBetaJpy & \dcAbsJpy & \dcAbsTJpy \\
    Canadian dollar & \dcSlCad & \dcBetaCad & \dcAbsCad & \dcAbsTCad \\
    \bottomrule
  \end{tabular}
\end{table}

\FloatBarrier
% =====================================================================================
\section{Can a clone do the job?}
\label{sec9:replication}

If a fund's return is mostly a set of exposures to liquid markets, a portfolio of those markets with
the same exposures could deliver it at a lower fee. \textcite{hasanhodzic2007} fitted such linear
clones to many hedge funds and found that for several categories of fund a clone built from a handful
of liquid factors came close to the funds' own returns, while for others it did not. A clone is a regression: at the end
of each month, regress the fund's last 24 monthly returns on the returns of a set of instruments, with
no intercept, and hold the fitted coefficients for the next month. Two results describe what can go
wrong.

\begin{proposition}[The cost of estimating a clone]\label{prop9:clone}
Let $y=x^\top\beta+e$, where the $k$ regressors $x$ are $\N(0,\Sigma)$ and $e$ is $\N(0,\sigma^2)$ and
independent of them. Fit $\hat\beta$ by least squares, without an intercept, on $T>k+1$ independent
draws, and apply it to a new independent draw $(x_0,y_0)$. Then
\begin{equation}\label{eq9:clone}
  \E\bigl[(y_0-x_0^\top\hat\beta)^2\bigr] = \sigma^2\Bigl(1+\frac{k}{T-k-1}\Bigr),
\end{equation}
while the mean squared residual within the fitting window has expectation $\sigma^2(T-k)/T$.
\end{proposition}

\begin{proof}
With $X$ the $T\times k$ matrix of regressors, $\hat\beta-\beta=(X^\top X)^{-1}X^\top e$, so
$y_0-x_0^\top\hat\beta=e_0-x_0^\top(X^\top X)^{-1}X^\top e$. Given $X$ and $x_0$ the two terms are
independent with mean zero, and the variance is $\sigma^2\bigl(1+x_0^\top(X^\top X)^{-1}x_0\bigr)$. Since $x_0$
is independent of $X$, $\E[x_0^\top(X^\top X)^{-1}x_0]=\operatorname{tr}\bigl(\Sigma\,\E[(X^\top X)^{-1}]\bigr)$, and
$X^\top X$ is Wishart with $T$ degrees of freedom and scale $\Sigma$, whose inverse has mean
$\Sigma^{-1}/(T-k-1)$ \parencite{muirhead1982}. The trace is then $k/(T-k-1)$. Within the window the
residual sum of squares is $\sigma^2$ times a chi-squared variable with $T-k$ degrees of freedom.
\end{proof}

A clone on seven instruments fitted to 24 months has an expected squared error out of sample of
$1+7/16=\pvTheory$ times the irreducible variance; a simulation of \pvSims{} fits gives \pvSim{}
(standard error \pvSimSe). Within the window it looks better than it is, with a residual variance of
\pvInSample{} times the truth, so the clone's error out of sample is about twice what the fit suggests.

\begin{proposition}[No linear clone of a straddle]\label{prop9:linear}
Let $X$ have a distribution symmetric about zero with $0<\E X^2<\infty$, and let $g$ be an even function
with $\E g(X)^2<\infty$. Then $\Cov(X,g(X))=\E[Xg(X)]=0$, so in the least squares regression of $g(X)$ on a
constant and $X$, and also without the constant, the coefficient on $X$ is zero:
the best linear clone of the payoff $g(X)$ holds no position in $X$.
\end{proposition}

\begin{proof}
$Xg(X)$ is an odd function of $X$, so its mean is zero by symmetry, and $\E X=0$ for the same reason.
Hence $\Cov(X,g(X))=\E[Xg(X)]-\E X\,\E g(X)=0$, and the coefficient $\Cov(X,g(X))/\Var X$ is zero.
\end{proof}

By \cref{prop9:straddle} the payoff of the linear trend rule is an even function of the total move less
the realised variance, and sign rules like the programme's are convex in sustained moves without
having exactly that payoff. The
result rules out a static linear exposure to the total move as a replica of its symmetric, convex
part; it does not by itself measure the replication error of a dynamic trend strategy across many
markets. Capturing the programme's nonlinear exposure needs dynamic trading rules, such as positions
in other trend rules, or instruments with nonlinear payoffs.

\paragraph{A test on a hypothetical manager}
To see both results at work, take a hypothetical manager who runs the programme of
\cref{sec9:programme} with an equal blend of six and twelve month lookbacks and adds returns of its
own, independent normal with a volatility of \rpIdio\% a year. The manager's rules are unknown to the
clone. The first clone regresses the manager on the seven market returns; the second on three trend
factors, the programmes of \cref{sec9:programme} with lookbacks of one, three and twelve months. Both
are refitted every month on the last \rpWin{} months, which leaves \rpMonths{} months out of sample
(\cref{fig9:clone}).

\begin{figure}[!tb]
  \centering
  \includegraphics[width=0.8\textwidth]{ch09/fig07_clone}
  \caption{Growth of one unit in the excess return of a hypothetical trend manager and of two clones
  refitted monthly on the last \rpWin{} months, out of sample, log scale. The manager is simulated from
  the programme of \cref{sec9:programme} as described in the text.}
  \label{fig9:clone}
\end{figure}


Over the full sample the seven markets explain only \rpFullA{} of the manager's variance in a single
regression, as \cref{prop9:linear} leads one to expect. Within each window the market clone fits
better, with an average $R^2$ of \rpInA, but out of sample its $R^2$ is \rpOutA: it does worse than
holding nothing. Its correlation with the manager is \rpCorrA, its tracking error \rpTeA\% a year, and
its Sharpe ratio \rpSRa{} against the manager's \rpSRt. The clone on trend factors has an average window
$R^2$ of \rpInB{} and an out of sample $R^2$ of \rpOutB, a correlation of \rpCorrB, a tracking error of
\rpTeB\%, and a Sharpe ratio of \rpSRb; it trails the manager by \rpMeanGapB\% a year, against
\rpMeanGapA\% for the market clone.
\Cref{fig9:clonecoef} shows why the market clone fails. Its coefficients swing from one refit to the
next, by \cfChangeA{} a month on average against \cfChangeB{} for the clone on trend factors, and the
median market coefficient changes sign \cfFlipsMed{} times in \cfMonths{} months. That is the noise of
\cref{prop9:clone}: seven coefficients estimated from \rpWin{} months. The factor clone's weight on the
twelve month factor stays between \cfTwelveMin{} and \cfTwelveMax.

\begin{figure}[!tb]
  \centering
  \includegraphics[width=0.96\textwidth]{ch09/fig18_clonecoef}
  \caption{Coefficients of each monthly refit of the two clones over the previous \rpWin{} months, on the
  seven market returns (left) and on the three trend factors (right).}
  \label{fig9:clonecoef}
\end{figure}


The second clone succeeds partly because the factors were built with the same seven markets and the
same kind of rule as the manager, which is a generous setting. For a real fund the factors have to be
guessed, and \cref{prop9:clone} charges for every factor added. The practical conclusion is narrow.
Where a strategy is a known family of rules on liquid markets, its return can be largely bought
without the manager. Where the return comes from illiquid positions or from selling insurance, a
clone built only from linear exposures to liquid markets can match the calm months and miss the
nonlinear losses of liquidity stress or of insurance claims, which is where the fund's risk lies.

\FloatBarrier
% =====================================================================================
\section{Fee mechanics}
\label{sec9:fees}

A hedge fund charges a management fee, a fixed share of assets each year, and an incentive fee, a share
$f$ of gains. The details of the incentive fee decide how much of a fund's return the investor keeps.
\begin{itemize}
  \item A \emph{high-water mark} is the highest value per unit on which a fee has been charged; the
        incentive fee is charged only on value above it, so a fund that loses must recover the loss
        before it is paid again.
  \item A \emph{hurdle} $h$ is a return the fund must earn before the incentive fee applies. Under a
        \emph{hard} hurdle the fee is charged only on the return above $h$. Under a \emph{soft} hurdle with
        full \emph{catch-up}, once the return exceeds $h$ the manager takes all of the excess until the
        fee equals $f$ of the whole return, and $f$ of the return after that.
  \item \emph{Crystallisation} is the date at which an accrued incentive fee becomes payable, usually
        the end of each year.
\end{itemize}

\begin{proposition}[Hurdles]\label{prop9:hurdle}
For a period return $R$, a rate $f\in(0,1)$ and a hurdle $h\ge0$, the incentive fees with no hurdle, a
hard hurdle and a soft hurdle with full catch-up are $fR^+$, $f(R-h)^+$ and
\begin{equation}\label{eq9:soft}
  \min\bigl((R-h)^+,\,fR^+\bigr)=(R-h)^+-(1-f)\bigl(R-h/(1-f)\bigr)^+,
\end{equation}
so the catch-up ends at a return of $h/(1-f)$ and the three fees are ordered hard $\le$ soft $\le$ none. If
$R\sim\N(\mu,s^2)$ their means are $fG(\mu,s)$, $fG(\mu-h,s)$ and $G(\mu-h,s)-(1-f)\,G\bigl(\mu-h/(1-f),s\bigr)$, with
$G(\mu,s)=\E[(\mu+sZ)^+]=\mu\Phi(\mu/s)+s\varphi(\mu/s)$ for a standard normal $Z$, the function of
\cref{eq9:callmean}.
\end{proposition}

\begin{proof}
For $R\le h$ both sides of \cref{eq9:soft} are zero. For $h<R\le h/(1-f)$, $R-h\le fR$, so the minimum is $R-h$, and
the second term on the right vanishes. For $R>h/(1-f)$ the minimum is $fR$, and the right side is
$R-h-(1-f)R+h=fR$. The ordering follows from $(R-h)^+\le R^+$ and from \cref{eq9:soft}: the soft fee is at
least $f(R-h)^+$, since $R-h\ge f(R-h)$ and $fR\ge f(R-h)$ when $R\ge h$, and at most $fR^+$. The means follow from
$\E(R-c)^+=G(\mu-c,s)$.
\end{proof}

With $f=\feF\%$, $h=\feH\%$ and a gross return with mean \feMu\% and volatility \feSig\%, the expected fee is
\feNone\% of assets with no hurdle, \feSoft\% with a soft hurdle and \feHard\% with a hard one; a simulation of a
million years gives the same figures to two decimals. The soft hurdle sounds like protection and costs
the manager little, because the catch-up ends at \feCatch\%, below the mean return. The hard hurdle takes
a third off the fee (\cref{fig9:fees}, panel (a)).

\begin{figure}[!tb]
  \centering
  \includegraphics[width=\textwidth]{ch09/fig11_fees}
  \caption{(a) Incentive fee at \feF\% on a year's gross return with no hurdle, a hard hurdle of \feH\% and a
  soft hurdle of \feH\% with full catch-up, from \cref{prop9:hurdle}. (b) Value at the start of
  \feYears{} years of annual incentive fees at \feF\%, as a share of the starting assets, when the assets grow at
  the bill rate of \feRf\% under the pricing measure: charged on each year's gain, by simulation and by
  \cref{prop9:hwm}, and with a high-water mark, by simulation.}
  \label{fig9:fees}
\end{figure}

\paragraph{The fee as an option}
Seen from the start of the year, the incentive fee is a call option that the investor writes to the
manager on the fund's assets. \textcite{goetzmann2003} value the fees of a fund with a high-water mark
as claims on its assets in this spirit. For one year the value follows from the Black-Scholes
formula of Chapter~5.

\begin{proposition}[The value of one year's incentive fee]\label{prop9:feecall}
Let the fund's assets before the incentive fee grow at the rate $r$ under the pricing measure with
volatility $s$, from a value of one. The value at the start of a fee of $f$ on the gain above $K$ at time
$T$ is
\begin{equation}\label{eq9:feecall}
  c=f\bigl[\Phi(d_1)-Ke^{-rT}\Phi(d_1-s\sqrt T)\bigr],\qquad d_1=\frac{-\log K+rT}{s\sqrt T}+\frac{s\sqrt T}2 .
\end{equation}
With a hurdle at the bill rate, $K=e^{rT}$, this is $c=f\bigl(2\Phi(s\sqrt T/2)-1\bigr)\approx fs\sqrt{T}/\sqrt{2\pi}$.
This is a value under the pricing measure of a lognormal model; the expected fee under the fund's own
distribution of returns is a different number.
\end{proposition}

\begin{proof}
The fee is $f(V_T-K)^+$, and $f$ times the price of a call on $V$ with strike $K$ is \cref{eq9:feecall}. With
$K=e^{rT}$, $d_1=s\sqrt T/2$ and $d_1-s\sqrt T=-s\sqrt T/2$, so $c=f(\Phi(s\sqrt T/2)-\Phi(-s\sqrt T/2))$. Near zero
$2\Phi(x)-1\approx2x\varphi(0)=2x/\sqrt{2\pi}$, and $x=s\sqrt T/2$.
\end{proof}

For a fund with a volatility of \feOptVol\% and a bill rate of \feRf\%, a \feF\% fee on the whole gain is worth
\feOptAtm\% of assets a year, and with a hurdle at the bill rate \feOptFwd\%, which the approximation gives to the
same two decimals. With the hurdle at the bill rate the value is close to proportional to the fund's
volatility, and without it the value still rises with volatility. A manager paid this way gains
from taking more risk, whatever that risk does for the investor.

\paragraph{The high-water mark}
Over several years the fee depends on the path. A high-water mark has a simple accounting identity.

\begin{proposition}[Fees and the high-water mark]\label{prop9:hwm}
Let a fund's value per unit be $N_0=1$ and, at the end of each period, $P_t=N_{t-1}(1+R_t)$ before the
incentive fee, $F_t=f(P_t-H_{t-1})^+$ and $N_t=P_t-F_t$, where $H_t=\max(H_{t-1},N_t)$ and $H_0=1$. Then
for every $T$,
\begin{equation}\label{eq9:hwm}
  \sum_{t=1}^TF_t=\frac f{1-f}\,(H_T-H_0).
\end{equation}
The left side is the undiscounted sum of the fees, each in units of the starting value per unit,
and not their present value.
If instead the fee is charged each period on that period's gain, $F'_t=f(P'_t-N'_{t-1})^+$, and the
returns $R_t$ are independent with the assets growing at $r$ under the pricing measure, the value of
the fees over $T$ periods is $1-(1-c)^T$, with $c$ the value of one period's fee from \cref{eq9:feecall}
with $K=1$ and $T$ one period.
\end{proposition}

\begin{proof}
If $P_t\le H_{t-1}$ then $F_t=0$ and $N_t=P_t\le H_{t-1}$, so $H_t=H_{t-1}$. If $P_t>H_{t-1}$ then
$N_t-H_{t-1}=(1-f)(P_t-H_{t-1})>0$, so $H_t=N_t$ and $F_t=\frac f{1-f}(H_t-H_{t-1})$. In both cases
$F_t=\frac f{1-f}(H_t-H_{t-1})$, and the sum telescopes. For the second part, a period's fee is
$N'_{t-1}f(R_t)^+$, and by independence its value at the start of the period is $N'_{t-1}c$, so the
discounted value of what remains in the fund falls by the factor $1-c$ each period. The value of all
fees is the starting value less the value of what the investor is left with, $1-(1-c)^T$.
\end{proof}

\Cref{eq9:hwm} says that over any horizon the fees are fixed by how far the high-water mark has risen,
whatever the path; the code checks it on every simulated path to within \feIdGap. A fund that gains 20\%,
loses 25\% and gains 30\%, with $f=\feF\%$, ends with a high-water mark \feExMark\% above its start and pays
\feExFeeH\% in incentive fees in all, summed without discounting as a share of the starting value;
charged on each year's gain the same returns cost \feExFeeR\% on the same basis. These are accounting
totals. The second part of the proposition values the difference under the pricing measure. Over \feYears{} years at a volatility of \feTenVol\%,
annual fees on each year's gain are worth \feTenClosed\% of the starting assets by the formula and \feTenReset\%
in a simulation of \feSims{} paths (standard error \feTenResetSe{} points); with a high-water mark they are worth
\feTenHwm\%, \feTenRatio\% as much (\cref{fig9:fees}, panel (b)). The value of the high-water mark to the investor
grows with the fund's volatility, since a volatile fund more often has losses to recover.


\FloatBarrier
% =====================================================================================
\section{Funds of funds}
\label{sec9:fof}

A fund of funds holds a portfolio of hedge funds and charges a fee of its own, typically a management
fee of about 1\% and an incentive fee of about 10\% of gains. It offers diversification across
managers, due diligence and access to funds that are closed to small investors. Diversification
across managers works like diversification across markets: with $n$ funds of volatility $\sigma$ and
pairwise correlation $\varrho$, the equally weighted pool has volatility $\sigma\sqrt{\varrho+(1-\varrho)/n}$,
which falls quickly at first and then levels off at $\sigma\sqrt\varrho$. The fees behave differently.
Each underlying fund charges its incentive fee on its own gain, whether or not the pool as a whole
gained, a point made by \textcite{brown2003}.

\begin{proposition}[Incentive fees that do not net]\label{prop9:netting}
For any gains $x_1,\dots,x_n$ and weights $v_i\ge0$, $\sum_iv_ix_i^+\ge\bigl(\sum_iv_ix_i\bigr)^+$. If the gains are
jointly normal, each $\N(\mu,\sigma^2)$ with pairwise correlation $\varrho\ge-1/(n-1)$, and the weights are $1/n$, the
expected incentive fee at rate $f$ is $f\,G(\mu,\sigma)$ per unit of assets when each fund charges on its own
gain, and $f\,G(\mu,\sigma_n)$ when gains are netted, with $\sigma_n=\sigma\sqrt{\varrho+(1-\varrho)/n}$ and
\begin{equation}\label{eq9:callmean}
  G(\mu,s)=\E\bigl[(\mu+sZ)^+\bigr]=\mu\,\Phi(\mu/s)+s\,\varphi(\mu/s),
\end{equation}
with $G(\mu,0)=\max(\mu,0)$ by continuity, the value at the boundary $\varrho=-1/(n-1)$, where $\sigma_n=0$.
$G$ increases with $s$, and $\sigma_n$ falls with $n$ when $\varrho<1$, so the gap between the two fees grows with $n$.
\end{proposition}

\begin{proof}
The positive part is convex and positively homogeneous, hence subadditive:
$(a+b)^+\le a^++b^+$ and $(va)^+=va^+$ for $v\ge0$; induction gives the inequality. Each fund's fee has
mean $f\E[x_i^+]=fG(\mu,\sigma)$, and the weighted mean of these is the same. The pooled gain is normal with
mean $\mu$ and variance $\sigma_n^2$. For \cref{eq9:callmean}, $\E[(\mu+sZ)^+]=\int_{-\mu/s}^\infty(\mu+sz)\varphi(z)\dd z
=\mu\Phi(\mu/s)+s\varphi(\mu/s)$, and differentiating in $s$ gives $\partial G/\partial s=\varphi(\mu/s)>0$.
\end{proof}

Take funds with expected annual gains of \ffMu\% before incentive fees, a volatility of \ffSig\% and
pairwise correlation \ffRho, each charging \ffPerf\% of its gain. To isolate the incentive fee, the
example leaves out the funds' management fees. Charged fund by fund, the expected incentive fee is
\ffSepOne\% of assets for any number of funds, \ffRateSep\% of the expected gain. If gains were netted
across the pool it would be \ffPoolFive\% with five funds, \ffPoolTen\% with ten and \ffPoolLimit\% in
the limit, \ffRateLimit\% of the gain. A simulation of ten funds gives \ffSimSep\% and \ffSimPool\% for
the two (\cref{fig9:fof}). In the years when the pool as a whole lost money, \ffLossYears\% of them, it
paid incentive fees in almost every case (\ffLossYearsPaid\% of all years), because some fund in the pool
had gained.

\begin{figure}[!tb]
  \centering
  \includegraphics[width=\textwidth]{ch09/fig08_fof}
  \caption{(a) Volatility of an equally weighted pool of $n$ funds with volatility \ffSig\% and pairwise
  correlation \ffRho, levelling off towards the limit $\sigma\sqrt\varrho$, the dashed line. (b) Expected incentive fees at
  \ffPerf\% when each fund charges on its own gain and when gains are netted, from \cref{eq9:callmean}.}
  \label{fig9:fof}
\end{figure}


The second layer adds to this. With a fund of funds fee of \ffFmg\% of assets and \ffFpf\% of its own
net gain, charged after the funds' incentive fees and with the funds' own management fees left out,
the investor's expected return in this one year simulation of ten funds falls from \ffMu\% before
incentive fees to \ffNetFof\%: fees take
\ffAllFees\% a year, \ffFeeShare\% of the expected gain. Ten funds cut the volatility from \ffVolOne\% to
\ffVolTen\%, and most of that cut is available with five (\ffVolFive\%). A fund of funds therefore
sells diversification that the arithmetic says levels off early, at a price that does not.

\FloatBarrier
% =====================================================================================
\section{Digital assets}
\label{sec9:crypto}

Bitcoin, proposed by \textcite{nakamoto2008}, is a digital token whose ownership is recorded on a
public ledger maintained by a network of computers, with a supply fixed by its protocol. Ether is the
token of a network that also runs programs. Neither gives its holder a contractual claim on cash
flows, so neither has a value that follows from discounting promised income; staking rewards and the
use of the network are benefits of another kind and need separate treatment. Their prices are set by what buyers will pay, and the questions for
an allocator are statistical: how risky they are, how they move with other assets, and how much a
short history can say.

\paragraph{How the ledger works}
A few facts about the protocol matter for an investor, and they can be stated without technical
detail. Transactions are grouped into blocks, and each block contains a cryptographic hash of the one
before, so the blocks form a chain in which changing an old transaction would change every later
block. A new block is accepted only with a \emph{proof of work}: a number that makes the block's hash
fall below a target, which can be found only by trial. Whoever finds it is paid in new coins, the
\emph{block subsidy}, plus the fees of the transactions in the block. The protocol adjusts the target
every 2016 blocks to keep the average time between blocks near ten minutes, and halves the subsidy
every 210,000 blocks, about every \ocEpochYears{} years at that spacing. In the reference software,
Bitcoin Core \parencite{bitcoincore}, the subsidy starts at 50 coins and is halved by a binary shift
of its value in the smallest unit, $10^{-8}$ of a coin.

\begin{proposition}[The supply limit]\label{prop9:supply}
With a subsidy of $\lfloor5\cdot10^9/2^j\rfloor$ units for each of the 210,000 blocks of period
$j=0,1,\dots$, the total number of coins ever issued is less than 21 million, and the first $J$ periods
issue $21\cdot10^6(1-2^{-J})$ coins up to the rounding of the subsidy.
\end{proposition}

\begin{proof}
Without rounding, period $j$ issues $210{,}000\times50\times2^{-j}=10.5\cdot10^6\cdot2^{-j}$ coins, and the
geometric series gives $21\cdot10^6(1-2^{-J})$ for the first $J$ periods and $21\cdot10^6$ in the limit.
Rounding down only removes units, and from $j=33$ on the subsidy is $\lfloor5\cdot10^9/2^{j}\rfloor=0$, so the total
is a finite sum strictly below the limit.
\end{proof}

The code sums the schedule exactly: \ocSupply{} coins, \ocShortfall{} short of 21 million, with the subsidy
reaching zero after \ocHalvings{} halvings. Half of all coins are issued in the first period and
\ocShareFour\% by the end of the fourth (\cref{fig9:onchain}, panel (a)). The supply is therefore known in
advance, and the subsidy that pays for the network's security falls by half every four years or so;
neither fact says anything about the price.
An Octave cross-check in \texttt{code/ch09/snippets/issuance.m} repeats the sum in floating point;
every partial sum is an integer below $2^{53}$, so it agrees exactly.

\begin{figure}[!tb]
  \centering
  \includegraphics[width=\textwidth]{ch09/fig14_onchain}
  \caption{(a) Bitcoins issued under the subsidy schedule of \cref{prop9:supply}, by period of 210,000
  blocks; the dashed line is 21 million. (b) Annualised volatility of daily log returns within each
  calendar year: bitcoin and ether from Coinbase prices, every day of the week, and US equities from the
  French daily market return, trading days only; 2026 runs to the end of September for the coins and
  August for equities.}
  \label{fig9:onchain}
\end{figure}


The second fact is how hard it is to rewrite the chain. \textcite{nakamoto2008} models an attacker
who controls a share $q$ of the computing power and tries to build a longer chain in secret, starting
$z$ blocks behind the honest one, whose share is $p=1-q$.

\begin{proposition}[Catching up]\label{prop9:ruin}
Let each new block be found by the honest miners with probability $p$ and by the attacker with
probability $q=1-p<p$, independently. The probability that the attacker, $z$ blocks behind, ever draws
level is $(q/p)^z$.
\end{proposition}

\begin{proof}
Let $D_n$ be the deficit after $n$ blocks, a random walk that rises by one with probability $p$ and
falls by one with probability $q$, from $D_0=z$. Then $(q/p)^{D_n}$ is a martingale, since
$p(q/p)^{d+1}+q(q/p)^{d-1}=(q/p)^d(q+p)$. Stop it when $D$ first reaches $0$ or some $M>z$. The stopped
martingale is bounded, so optional stopping gives
$(q/p)^z=\Pr(\text{0 first})+(q/p)^M(1-\Pr(\text{0 first}))$, and
$\Pr(\text{0 first})=\bigl((q/p)^z-(q/p)^M\bigr)/\bigl(1-(q/p)^M\bigr)$. Let $M\to\infty$; since $q/p<1$ the
limit is $(q/p)^z$.
\end{proof}

Nakamoto's own calculation allows for the blocks the attacker has already found while the honest chain
built its $z$, treating that number as Poisson with mean $zq/p$ and weighting the ruin probability by
it. The code reproduces the paper's table for $q=0.1$ to all seven decimals given there, and the number
of blocks needed to bring the probability below 0.1\% for all \ocZCount{} values of $q$ in the paper's
second table, from 5 blocks for $q=0.1$ to \ocZLast{} for $q=\ocQLast$. With $q=0.1$ and six blocks the ruin probability is \ocRuinSix{}
and Nakamoto's \ocNakSix. The arithmetic assumes the attacker's share is below a half and fixed; it says
nothing about other risks to the holder, such as the loss of keys or the failure of an exchange.

\paragraph{The record}
From \btFirst{} to \btLast, \btMonths{} months, the Coinbase price of bitcoin rose from \btPriceStart{} to
\btPriceEnd{} US dollars, a compound return of \btGeo\% a year. Its monthly returns had an annualised
volatility of \btVol\%, a worst month of \btWorst\% and a best of \btBest\%, with skewness \btSkew{} and
excess kurtosis \btKurt. On daily prices its worst fall from a peak was \btMddDaily\%, and its longest
spell below a previous peak lasted \btUnderYears{} years. Ether, from \ceFirst{} (\ceMonths{} months),
compounded at \ceGeo\% a year with a volatility of \ceVol\%, a worst month of \ceWorst\%, a best of
\ceBest\% and an excess kurtosis of \ceKurt; its worst fall was \ceMddDaily\% and its longest spell
under water \ceUnderYears{} years. The two had a monthly correlation of \btCorrEth{}
(\cref{fig9:crypto}, panel (a)).

\begin{figure}[!tb]
  \centering
  \includegraphics[width=\textwidth]{ch09/fig09_crypto}
  \caption{(a) Drawdowns of bitcoin and ether from their previous peaks, daily Coinbase prices,
  December 2014 (ether May 2016) to September 2026. (b) Correlation of monthly bitcoin returns with US
  equity excess returns over trailing 24 month windows, to \tfLast.}
  \label{fig9:crypto}
\end{figure}


\paragraph{What the record cannot say}
The average monthly return of bitcoin, annualised, was \btArith\% with a standard error of \btMeanSe\%.
The interval is wide, but the larger problem is that bitcoin is studied because it survived and grew.
At the start of the sample it was one of many tokens, and many of the others have since lost most of
their value or stopped trading; choosing the one that did best after the fact adds winner selection to
the survivorship bias of \cref{sec9:biases}, its strongest form. \Cref{exr9:ten} puts a number on it: among twenty tokens with no expected return and a
volatility of 100\% a year, the best ten year average is expected to be about \exJBest\% a year. A record
like bitcoin's says that the asset was very volatile and that its holders were rewarded after the
fact; it says little about the expected return of the next decade. \textcite{liu2021} find that
cryptocurrency returns are explained by factors specific to the market, momentum among them, and not by
exposures to stocks, currencies or commodities, in a sample that ends before the period of higher
correlation below.

\paragraph{Correlation}
From 2015 to 2019 the monthly correlation of bitcoin with US equity excess returns was \btCorrEarly. From
2020 to \tfLast{} it was \btCorrLate, with a beta of \btBetaLate{} on the equity market
(\cref{fig9:crypto}, panel (b)). Over the second period bitcoin had a large positive exposure to equities,
which cut its value as a diversifier, though its large risk unrelated to equities kept it far from a
leveraged equity position, and an allocator who relied on the early correlation would have
underestimated its contribution to portfolio risk.

\paragraph{Has the correlation changed?}
A change in correlation estimated from a few years of monthly data needs a test. The usual one uses
Fisher's transformation $z(\hat\varrho)=\operatorname{artanh}\hat\varrho=\tfrac12\log\bigl((1+\hat\varrho)/(1-\hat\varrho)\bigr)$.

\begin{proposition}[A test for a change in correlation]\label{prop9:fisher}
Let $\hat\varrho_n$ be the sample correlation of $n$ independent draws from a bivariate normal law with
correlation $\varrho\in(-1,1)$. Then $\sqrt n\bigl(\operatorname{artanh}\hat\varrho_n-\operatorname{artanh}\varrho\bigr)$ converges
in law to $\N(0,1)$. For two independent samples of sizes $n_1$ and $n_2$ with the same correlation,
\begin{equation}\label{eq9:fisher}
  Z=\frac{\operatorname{artanh}\hat\varrho_2-\operatorname{artanh}\hat\varrho_1}{\sqrt{1/(n_1-3)+1/(n_2-3)}}
\end{equation}
is therefore approximately standard normal.
\end{proposition}

\begin{proof}
For bivariate normal data $\sqrt n(\hat\varrho_n-\varrho)$ converges in law to $\N\bigl(0,(1-\varrho^2)^2\bigr)$, a standard
consequence of the central limit theorem for the sample moments. The derivative of $\operatorname{artanh}$ is
$1/(1-\varrho^2)$, so by the delta method the transformed estimate has asymptotic variance $1/n$, whatever
$\varrho$ is. The two estimates are independent, so their difference has variance close to $1/n_1+1/n_2$;
Fisher's refinement $1/(n-3)$ improves the approximation in small samples and has the same limit.
\end{proof}

With \ocNEarly{} months from 2015 to 2019 and \ocNLate{} from 2020 to \tfLast, \cref{eq9:fisher} gives $Z=\ocFisherZ$
and a two-sided $p$-value of \ocFisherP. The 95\% interval for the correlation is \ocEarlyLo{} to \ocEarlyHi{}
in the first period and \ocLateLo{} to \ocLateHi{} in the second. The rise is larger than sampling error
would explain on the test's terms, but the $p$-value is nominal: it assumes independent, bivariate
normal months and a break date fixed in advance, and it ignores how much bitcoin's own volatility
changed over the sample (next paragraph). A test chosen after looking at a chart also overstates its own significance,
since the break date was picked where the change looked largest.

The volatility has moved too. On daily prices, annualised, bitcoin's volatility within a calendar year
fell nearly threefold over the sample, from \btVolFirstYear\% in \btVolFirstYearName{} to \btVolLastYear\% in
\btVolLastYearName{}~(\cref{fig9:onchain}, panel (b)). Against US equities in the same year it
ranged from \ocRatioMin{} to \ocRatioMax{} times, and from 2022 to 2025 from \ocRatioRecentMin{} to
\ocRatioRecentMax{} times, the \ocRatioRecentMax{} in \ocRatioRecentMaxYear. A risk budget set
on one year's estimate could have been wrong by close to a factor of two a year later, as from 2016
to 2017.


\paragraph{How much risk a small position carries}
The \emph{risk contribution} of a holding is its weight times the covariance of its return with the
portfolio's, and the risk contributions add up to the portfolio's variance (Chapter~7). For a sleeve of
weight $w$ in a portfolio otherwise invested in $P$, the share of risk has a closed form.

\begin{proposition}[Weight for a given risk share]\label{prop9:share}
Let a portfolio hold $1-w$ in $P$ and $w$ in $B$, with $\Var P=a>0$, $\Var B=b>0$, $\Cov(P,B)=c$ and
$c^2<ab$. The share of the portfolio's variance due to $B$ is
\begin{equation}\label{eq9:share}
  s(w) = \frac{w^2b+w(1-w)c}{(1-w)^2a+2w(1-w)c+w^2b}.
\end{equation}
For each $s\in(0,1)$ there is exactly one $w\in(0,1)$ with $s(w)=s$, namely $w=u/(1+u)$ with $u$ the
positive root of $(1-s)b\,u^2+(1-2s)c\,u-s\,a=0$. If $c\ge0$, $s(w)$ increases with $w$.
\end{proposition}

\begin{proof}
The contribution of $B$ is $w\Cov(B,(1-w)P+wB)=w^2b+w(1-w)c$, and the denominator of \cref{eq9:share} is
the portfolio variance, positive since $c^2<ab$. Put $u=w/(1-w)$, a bijection from $(0,1)$ to
$(0,\infty)$, and divide numerator and denominator by $(1-w)^2$: $s=(u^2b+uc)/(a+2uc+u^2b)$. Clearing the
denominator gives the quadratic, whose leading coefficient $(1-s)b$ is positive and whose constant
$-sa$ is negative, so it has exactly one positive root. The derivative of $(u^2b+uc)/(a+2uc+u^2b)$ in $u$
has the sign of $2uab+ac+u^2bc$, which is positive when $c\ge0$.
\end{proof}

Over January 2015 to \tfLast{} a 60/40 portfolio of US equities and the US 10 year bond had a volatility
of \btVolSixty\%, and its correlation with bitcoin's excess return was \btCovSign. A weight of 1\% in
bitcoin would have carried \btShareOne\% of the portfolio's risk, 2\% would have carried \btShareTwo\% and 5\%
would have carried \btShareFive\%. By \cref{prop9:share}, bitcoin carries a tenth of the risk at a weight of
\btWTen\%, a quarter at \btWQuarter\% and half at \btWHalf\%. These are estimates from one volatile decade,
and they move with the estimates. \Cref{fig9:riskshare} redraws the risk share with each half of the
period on its own. Over 2015 to 2019 the correlation was \rsCorrEarly, and a tenth of the risk came at a
weight of \rsWTenEarly\% and half at \rsWHalfEarly\%; from 2020 the correlation was \rsCorrLate, and the
same shares needed \rsWTenLate\% and \rsWHalfLate\%. Above a weight of \rsCross\% the later estimates
give bitcoin the smaller share despite the higher correlation, because bitcoin's volatility fell from
\rsVolBtcEarly\% to \rsVolBtcLate\% a year while that of the 60/40 portfolio rose from
\rsVolSixtyEarly\% to \rsVolSixtyLate\%. A risk budget set on one period's estimates can be off by a
factor of two in the weight.

\begin{figure}[!tb]
  \centering
  \includegraphics[width=0.74\textwidth]{ch09/fig19_riskshare}
  \caption{Bitcoin's share of the variance of a 60/40 portfolio against its weight, by \cref{prop9:share},
  with the variances and covariance estimated over the whole period and over each half. The dotted lines
  mark shares of 10\%, 25\% and 50\%.}
  \label{fig9:riskshare}
\end{figure}


\FloatBarrier
% =====================================================================================
\section{Case study: a trend fund and a bitcoin sleeve in a 60/40 portfolio}
\label{sec9:case}

The case study asks what the two additions of this chapter did to a portfolio that holds 60\% in US
equities and 40\% in the US 10 year bond, rebalanced monthly. The first addition is the trend programme
of \cref{sec9:programme} after the 2 and 20 fee, held as a fund at 10\% of the portfolio, from
\tfFirst{} to \tfLast. The second is a bitcoin position of 2\% or 5\%, from \btFirst, the first month
with a return from the Coinbase series, which starts on \btSeriesStart, to \tfLast. Each addition is funded pro rata from the 60/40 holdings. Both
are frictionless backtests at month-end prices, and the trend fund is a constructed programme, so the
results describe what these rules and weights would have done, not what an investor in any actual fund
received. \Cref{tab9:case} gives the results, and \cref{fig9:case} the trend case.

\begin{figure}[!tb]
  \centering
  \includegraphics[width=\textwidth]{ch09/fig10_case}
  \caption{(a) Growth of one unit in the 60/40 portfolio and in the portfolio with 10\% in the trend
  fund after fees, \tfFirst{} to \tfLast, log scale. (b) Their drawdowns.}
  \label{fig9:case}
\end{figure}

\begin{table}[t]
  \centering
  \small
  \caption{Annual excess return over bills, volatility, Sharpe ratio, worst drawdown and worst month of
  the 60/40 portfolio and of the portfolios with each addition, all from monthly excess returns over
  bills; drawdowns compound them. Returns are in percent, a year for the mean. The change in the
  Sharpe ratio against the 60/40 portfolio of the same period is followed by a 95\% circular block
  bootstrap interval, blocks of \csBlock{} months and \csReps{} samples; changes are computed before
  rounding.}
  \label{tab9:case}
  \setlength{\tabcolsep}{2.5pt}
  \begin{tabular}{@{}lrrrrrl@{}}
    \toprule
    & Excess & Volatility & Sharpe & Drawdown & Worst month & Change in Sharpe\\
    \midrule
    \multicolumn{7}{@{}l}{\tfFirst{} to \tfLast}\\
    60/40 & \csAMean & \csAVol & \csASR & \csAMdd & \csAWorst & \\
    with 10\% trend fund, net & \csBMean & \csBVol & \csBSR & \csBMdd & \csBWorst & \csDiffTab{} [\csDiffLo, \csDiffHi]\\
    with 10\% trend, before fees & \csGMean & \csGVol & \csGSR & \csGMdd & \csGWorst & \csDiffG{} [\csDiffGLo, \csDiffGHi]\\
    \addlinespace
    \multicolumn{7}{@{}l}{\btFirst{} to \tfLast}\\
    60/40 & \ccAMean & \ccAVol & \ccASR & \ccAMdd & \ccAWorst & \\
    with 2\% bitcoin & \ccTwoMean & \ccTwoVol & \ccTwoSR & \ccTwoMdd & \ccTwoWorst & \ccTwoDiff{} [\ccTwoLo, \ccTwoHi]\\
    with 5\% bitcoin & \ccFiveMean & \ccFiveVol & \ccFiveSR & \ccFiveMdd & \ccFiveWorst & \ccFiveDiff{} [\ccFiveLo, \ccFiveHi]\\
    \bottomrule
  \end{tabular}
\end{table}


\paragraph{The trend fund}
The trend fund's correlation with the 60/40 portfolio was \csCorrAT. Adding it at 10\% lowered the
volatility from \csAVol\% to \csBVol\%, the worst drawdown from \csAMdd\% to \csBMdd\% and the worst month from
\csAWorst\% to \csBWorst\%. From November 2007 to February 2009 the portfolio with the fund lost \csBGfcAbs\%
against \csAGfcAbs\%, and in 2022 \csBTwoAbs\% against \csATwoAbs\%, both in total return, as in \cref{sec9:programme}, while
the drawdowns above are of excess returns. The cost was in the return: \csBMean\% a year
over bills against \csAMean\%, because the fund's net return was low. The Sharpe ratio rose by
\csDiff, from \csASR{} to \csBSR, and the bootstrap interval for the change, from \csDiffLo{} to
\csDiffHi, includes zero; in \csDiffP\% of the bootstrap samples the change was zero or negative. Before
fees the same allocation would have raised the Sharpe ratio to \csGSR. The trend fund bought a
smaller drawdown in the two worst episodes for this portfolio, and the fee took most of what it would
otherwise have added to the return.
\Cref{fig9:wealthgap} shows when the fund paid. The log wealth of the portfolio with the fund gained
\cwGfc{} points on 60/40 from the end of October 2007 to the end of February 2009, reaching a lead of
\cwPeak{} points in \cwPeakWhen, then gave back almost \cwFromPeakUp{} points to end \cwEndAbs{} points behind.
Before fees it would have ended \cwEndGrossAbs{} points behind, so the fee accounts for \cwFeeWedge{}
points. The Sharpe ratio still rose, because the volatility fell proportionally more than the return.

\begin{figure}[!tb]
  \centering
  \includegraphics[width=0.86\textwidth]{ch09/fig20_wealthgap}
  \caption{Log wealth of the 60/40 portfolio with 10\% in the trend fund less that of 60/40, in percentage
  points, net of the 2 and 20 fee and before fees, \tfFirst{} to \tfLast.}
  \label{fig9:wealthgap}
\end{figure}


\paragraph{The bitcoin sleeve}
From January 2015 the 60/40 portfolio earned \ccAMean\% a year over bills with a Sharpe ratio of \ccASR.
With 2\% in bitcoin the figures are \ccTwoMean\% and \ccTwoSR, and with 5\% they are \ccFiveMean\% and
\ccFiveSR, with worst drawdowns of \ccTwoMdd\% and \ccFiveMdd\% against \ccAMdd\%. Both bootstrap intervals
for the change in the Sharpe ratio lie above zero. In 2022, when bitcoin fell \ccBtcTwoAbs\%, the portfolio
with 2\% lost \ccTwoTwoAbs\% against \ccATwoAbs\% without it. That gain cannot be taken at face value. The bootstrap resamples months of the one asset that, among many candidates, rose from
\btPriceStart{} to \btPriceEnd{} dollars, so it measures the sampling error of this history and not the
selection that chose it, nor the choice of model, and \cref{exr9:ten} shows how large that selection can be. A more defensible
reading is the risk budget of \cref{prop9:share}: at 2\% bitcoin carried about \btShareTwo\% of the
portfolio's risk, an amount an investor can decide to accept or not without a forecast of its return.

\begin{figure}[!tbp]
  \centering
  \includegraphics[width=0.66\textwidth]{ch09/fig21_selection}
  \caption{Change in the Sharpe ratio of 60/40 from adding at 2\% the best ex post of \sxTokens{}
  simulated tokens with no expected excess return, in \sxReps{} simulations, January 2015 to \tfLast.
  The red line is the change from bitcoin, \sxActualFine{} before rounding.}
  \label{fig9:selection}
\end{figure}

\Cref{fig9:selection} puts selection in the units of the table. Draw \sxTokens{} tokens with bitcoin's
volatility and correlation with 60/40 but no expected excess return, conditional on the actual 60/40
months, add the best of them ex post at 2\%, and repeat \sxReps{} times. The best of twenty raises the
Sharpe ratio by \sxMedian{} in the median simulation and by \sxQ{} at the 95th percentile, and a gain as
large as bitcoin's \ccTwoDiff{} came up in \sxCount{} of the \sxReps{} simulations. In this particular
zero-mean Gaussian experiment, selection among twenty candidates rarely generates a gain as large as
bitcoin's; the result does not establish how much selection contributed in the actual, larger and
non-Gaussian cryptocurrency universe. From a separate run of \sxSingle{} single tokens, it would take
about \sxNeeded{} candidates for the median best to match it. Whether that is many depends on how many
candidates an investor in 2015 would really have weighed, which no backtest can say.



\FloatBarrier
% =====================================================================================
\section{Limitations}
\label{sec9:limits}

The trend programme is a sketch of a professional one. It trades seven markets where commercial
programmes trade many more, builds bond returns from the 10 year yield alone, without the roll-down
of \cref{sec9:carry}, omits the carry of currencies and the roll of oil, and pays a flat cost per unit
traded. Its signal and its risk target were chosen in line
with the published literature, which saw some of the same data, so even the gross record is not free of
hindsight.

The database and replication sections use simulated funds. The simulation shows how biases and
clone errors arise and how large they can be under stated rules; real databases have their own rules
for entry and exit, and the size of each bias there has to be measured from the database itself. The
fund of funds example assumes normal gains and a single year; over several years, with high-water
marks, each fund's fee also depends on its own path, as \cref{prop9:hwm} shows. The closure model has
one loss limit for every fund and no skill, and the fee values assume that the fund's volatility is
constant and known. The carry results rest on one currency pair and one bond market.

The digital asset record is short, comes from one exchange's prices, and covers an asset chosen after
the fact. The protocol arithmetic describes the rules as written in the reference software and
assumes an attacker's share is fixed. The correlations, the volatility and therefore the risk shares are estimates from a decade in
which the market's structure changed, and the next decade's figures may differ.

\FloatBarrier
% =====================================================================================
\section{Summary}
\label{sec9:summary}

\begin{itemize}
  \item Hedge fund strategies are best described by the risk they carry and the shape of their payoff.
        Many sell insurance against crashes or illiquidity; trend followers buy convexity.
  \item A trend rule on a martingale earns nothing before costs. What it earns splits into a reward
        for the drift and a timing part, and the linear trend rule has an exact payoff identity, one half of
        the squared total move less one half of the realised sum of squared returns, which gives it
        straddle-like convexity without making it a purchased straddle; the sign rules used in practice share its convexity only
        approximately.
  \item A seven market programme from 2001 earned a Sharpe ratio of \tfSR{} before fees, with a
        negative correlation to equities and gains in the long equity falls of 2008 and 2022. A 2 and
        20 fee cut the ratio to \tfSRnet.
  \item The Sharpe ratio of an equal mix of $n$ strategies with common correlation $\varrho$ is that
        of one strategy times $\sqrt{n/(1+(n-1)\varrho)}$.
  \item The simple return of a futures position is the carry plus the spot return times one plus the
        carry; in logs, under the cost of carry, it is the spot return plus carry. Carry was small beside
        the monthly risk of the Canadian dollar and the Canada 10 year bond. Over twenty five years carry
        alone gave about \gcCarryShare\% of the bond's excess return, and carry and roll-down \gcExpShare\%.
  \item Survivorship and backfill biases have closed forms under normal returns, and simple reporting
        rules produce biases of the order of a percent a year. Closure at a loss limit gives a bias
        that grows with the horizon, with a closed form from the reflection principle.
  \item A regression splits a fund's mean return exactly into alpha and factor parts, but the split
        depends on the factors chosen.
  \item A regression clone pays $1+k/(T-k-1)$ in squared error out of sample, and under a symmetric
        distribution the least squares slope of a payoff that is even in the move is zero, so no static
        linear clone can reproduce it.
  \item Incentive fees summed over any horizon, without discounting, are $f/(1-f)$ times the rise in the
        high-water mark. One year's fee is valued as a call on the fund's assets under the pricing
        measure, worth about $fs\sqrt T/\sqrt{2\pi}$ with a hurdle at the bill rate,
        and a high-water mark cuts the value of ten years of fees by about a third at a volatility of
        15\%.
  \item Incentive fees charged fund by fund exceed fees on the pooled gain, by more as the pool grows.
  \item Bitcoin's supply is fixed in advance below 21 million coins, and with fixed shares of computing power
        the chance that an attacker $z$ blocks behind catches up falls geometrically in $z$. Its record is volatile, short and selected. Its correlation with
        equities rose after 2019 by more than sampling error would explain on the nominal Fisher test,
        though the break date was chosen after the fact, and a 2\% position carried about
        \btShareTwo\% of a 60/40 portfolio's risk.
\end{itemize}

\FloatBarrier
% =====================================================================================
\section{Exercises}
\label{sec9:exercises}

\begin{enumerate}[label=\textbf{9.\arabic*}, leftmargin=*, ref=9.\arabic*]
  \item\label{exr9:one} A market returns 2\%, $-1$\%, 3\%, $-2$\% and 4\% in five periods. Compute the profit of
  the linear trend rule of \cref{prop9:straddle} period by period, and check it against
  \cref{eq9:straddle}.
  \item\label{exr9:two} The correlation between a market's past year return and its next month return
  is 0.05, and its monthly volatility is 4\%. Under the assumptions of \cref{prop9:sign}, what does the
  sign rule earn per month, and what is its annual Sharpe ratio?
  \item\label{exr9:three} Twenty trend sleeves each have a Sharpe ratio of 0.2 and pairwise correlation
  0.1. Find the Sharpe ratio of their equal mix and its limit as the number of sleeves grows.
  \item\label{exr9:four} Funds have independent normal monthly returns with mean 0.5\% and volatility 4\%.
  A database keeps only funds whose two year average return is positive. What share survives, and by
  how much does the survivors' average overstate the true mean, in percent a year?
  \item\label{exr9:five} Funds with monthly returns of mean 0.5\% and volatility 3\% join a database after
  a twelve month incubation with a positive average, bringing that history, and then report for four
  years. Find the mean of the incubation months and the bias of the full record, in percent a year.
  \item\label{exr9:six} A clone on ten factors is fitted to 36 months. By \cref{prop9:clone}, by what factor
  does its expected squared error out of sample exceed the irreducible variance? How long must the
  window be for the excess to fall to 10\%?
  \item\label{exr9:seven} A fund of funds holds two funds in equal amounts. In one year the first gains 30\%
  and the second loses 10\%, before fees, and each charges 20\% of its gain. What incentive fee does the
  fund of funds pay, as a share of its assets, what would it pay on the pooled gain, and what share of
  the pooled gain does it pay in fact?
  \item\label{exr9:eight} A portfolio has a volatility of 10\%. An asset with a volatility of 75\% has a
  correlation of 0.3 with it. At what weight does the asset carry a quarter of the risk of the
  combination, and what share does it carry at a weight of 3\%?
  \item\label{exr9:nine} A trend fund earns 9\% before fees in a year when bills pay 3\%. It charges 2\% of
  its starting assets and then 20\% of the remaining gain, with no hurdle and a high-water mark at the
  starting value. Find the
  investor's return and excess return, and the share of the gross excess return the investor keeps.
  \item\label{exr9:ten} Twenty tokens have independent normal returns with no expected return and a
  volatility of 100\% a year, independent from year to year. Ignore delisting and any selection before
  the end of the ten years. After ten years an analyst studies the one with the highest average
  return. What is its expected annual average? Use $\E\max(Z_1,\dots,Z_{20})$ for twenty independent
  standard normals, computed numerically.
  \item\label{exr9:eleven} A fund with a high-water mark charges 20\% of gains at the end of each year and
  has gross returns of 25\%, $-20$\% and 30\% in three years. Find the fee each year and the total, check
  the total against \cref{eq9:hwm}, and compare it with the fees charged on each year's gain.
  \item\label{exr9:twelve} With $f=20\%$ and $h=5\%$, find the incentive fee with no hurdle, a hard hurdle and a soft
  hurdle with full catch-up when the gross return is 10\%, and again when it is 6\%. Where does the
  catch-up end?
  \item\label{exr9:thirteen} Over one year a Canadian bill pays 3\% and a US bill 4.5\%, and the price of a US dollar falls
  from 1.35 to 1.33 Canadian dollars. Split the US dollar excess return of Canadian bills funded at the
  US rate into its spot, carry and cross parts, using \cref{eq9:fxcarry} with the annual rates as
  simple rates for the year in place of $i/12$.
  \item\label{exr9:fourteen} An attacker controls 20\% of the computing power and the honest chain is six blocks
  ahead. Find the probability of \cref{prop9:ruin} that the attacker ever draws level, and explain why
  Nakamoto's figure, \exNNak, is larger.
\end{enumerate}

\subsection*{Answers}
\begin{enumerate}[label=\textbf{9.\arabic*}, leftmargin=*]
  \item The cumulative moves before each period are 0, 2, 1, 4 and 2\%, so the profits are $0$, $-2$, $3$,
  $-8$ and $8$ in units of $10^{-4}$, a total of \exAPnl. The total move is \exAS\% and the sum of squared
  returns is \exAQ{} in the same units, so \cref{eq9:straddle} gives $(36-34)/2=\exAPnl$.
  \item $\sqrt{2/\pi}\times0.05\times4\%=\exBMonth\%$ a month. The position is one unit long or short, so the
  monthly volatility is 4\% and the annual Sharpe ratio is $\sqrt{12}\times0.16/4=\exBSR$.
  \item $0.2\sqrt{20/(1+19\times0.1)}=\exCSR$, and the limit is $0.2/\sqrt{0.1}=\exCLim$.
  \item The two year average has standard deviation $4\%/\sqrt{24}$, so $a=-0.5\sqrt{24}/4$ and
  $1-\Phi(a)=\exDShare\%$ survive. By \cref{eq9:survivor} the survivors average \exDMean\% a year against the
  true 6\%, a bias of \exDBias\%.
  \item By \cref{eq9:survivor} with $T=12$ the incubation months average \exEInc\% a year. By
  \cref{eq9:backfill} with $K=12$ and $M=48$ the full record overstates the mean by \exEBias\% a year.
  \item $1+10/25=\exFInf$. The excess $k/(T-k-1)$ is 0.1 when $T-11=100$, so $T=\exFT$ months.
  \item The first fund pays 20\% of 30\% on half the assets, \exGFee\% in all; the second pays nothing.
  On the pooled gain of 10\% the fee would be \exGNet\%. The fund of funds pays $3/10=\exGRate\%$ of its gain.
  \item With $a=0.01$, $b=0.5625$ and $c=0.3\times0.1\times0.75$, \cref{prop9:share} with $s=0.25$ gives a weight
  of \exHWx\%, about \exHW\%, and \cref{eq9:share} at $w=0.03$ gives a share of \exHSharex\%, about \exHShare\%.
  \item After the management fee of 2\% of the starting assets the gain is 7\%, and the incentive fee,
  with no hurdle and the high-water mark at the starting value, takes a fifth of it, leaving \exINet\%.
  The excess return is \exINetEx\%, \exIRatio\% of the gross excess return of 6\%.
  \item The ten year average of each token is normal with mean zero and standard deviation
  $100\%/\sqrt{10}$. $\E\max$ of twenty standard normals is \exJMax, so the best token's expected average
  is $\exJMax\times100/\sqrt{10}=\exJBest\%$ a year, though none of the twenty has any expected return.
  \item In the first year the value before the fee is 1.25, the fee \exKOne\% and the value 1.20, the new
  high-water mark. The second year ends at 0.96 with no fee. The third ends at 1.248 before the fee, which
  is \exKThree\% on the gain above 1.20. The sum of the three fees, undiscounted and as a share of the starting value, is \exKTot\%, and $0.2/0.8$ times the rise of the high-water
  mark, \exKRise\%, is the same. Charged on each year's gain the third year's fee is \exKResetThree\% and the total
  \exKReset\%.
  \item At 10\%: \exLNTen\%, \exLHTen\% and \exLSTen\%, since the catch-up has ended. At 6\%: \exLNSix\%, \exLHSix\% and
  \exLSSix\%, the whole return above the hurdle. The catch-up ends at $5/0.8=6.25\%$.
  \item The spot part is $1.35/1.33-1=\exMSpot\%$, the carry $3-4.5=\exMCarry\%$ and the cross term
  $0.03\times\exMSpot\%=\exMCross\%$, a total of \exMEx\%: the fall in the US dollar paid for the lower
  Canadian rate, with a little to spare.
  \item $(0.2/0.8)^6=\exNRuin$. The ruin probability assumes the attacker starts exactly six blocks behind.
  While the honest chain found its six blocks the attacker was mining too, and on average found
  $6\times0.25=1.5$ of its own, so it is usually fewer than six behind; averaging over that number raises the
  probability.
\end{enumerate}

\FloatBarrier
% =====================================================================================
\section{Notes, sources and credits}
\label{sec9:credits}

\paragraph{Sources}
Time series momentum is from \textcite{moskowitz2012}, and the century of evidence is in
\textcite{hurst2017}. The option-like returns of trend followers are from \textcite{fung2001}, those of
merger arbitrage from \textcite{mitchell2001} and those of hedge fund indices from
\textcite{agarwal2004}; the hidden market exposure of stale prices is from \textcite{asness2001}, with the
smoothing model of \textcite{getmansky2004}. The database biases were documented by
\textcite{brown1992} for mutual funds and by \textcite{fung2000}, \textcite{liang2000},
\textcite{ackermann1999} and \textcite{aiken2013} for hedge funds, why funds stop reporting is studied by
\textcite{grecu2007}, and the reporting biases are inferred from equity holdings by \textcite{agarwal2013}; the competition and survival of
hedge funds and commodity trading advisers are studied by \textcite{brown2001}. The reflection
principle behind \cref{prop9:closure} is in \textcite{shreve2004}, and the continuity correction for
monthly monitoring is from \textcite{broadie1997}. Carry across asset classes is from
\textcite{koijen2018}, roll returns in commodity futures from \textcite{erb2006}, and the deviations from
covered interest parity from \textcite{du2018}. Linear replication is
from \textcite{hasanhodzic2007}, and the mean of the inverse Wishart matrix used in \cref{prop9:clone}
is in \textcite{muirhead1982}. The fees of funds with high-water marks are valued as claims on the
fund by \textcite{goetzmann2003}, and the fees of funds of funds are analysed by \textcite{brown2003}.
Bitcoin was proposed by \textcite{nakamoto2008}, whose Section 11 gives the attacker calculation
reproduced in \cref{sec9:crypto}; the subsidy schedule and the supply limit are those of the reference
software, Bitcoin Core \parencite{bitcoincore}. The factor structure of cryptocurrency returns is
from \textcite{liu2021}. The straddle identity of \cref{prop9:straddle} is equation (7) of \textcite{dao2017}, who
relate trend following to the gap between long and short horizon variance. The sign rule, survivor,
backfill and closure formulas, the hurdle formulas, the
high-water mark identity, the netting inequality, the supply limit and the risk share weight are short
derivations given here for completeness.

\paragraph{Software used}
The scripts call NumPy \parencite{numpy} for arrays and random numbers, SciPy \parencite{scipy}
for the normal law and numerical integration, pandas \parencite{pandas} to read and sample the daily
files, and Matplotlib \parencite{matplotlib} for the figures. The cross-checks run in R
\parencite{rcore} and GNU Octave \parencite{octave}. \texttt{code/requirements-lock.txt} records the full
environment.

\paragraph{Projects for comparison}
The QuantStats package \parencite{quantstats} computes the Sharpe ratios, drawdowns and annual returns
of \cref{sec9:programme,sec9:case} from a return series and is a useful check on the code here. The full
list, with links, is in \texttt{CREDITS.md}.

\paragraph{Data}
US equity excess returns and the bill rate are from the data library of \textcite{frenchlib}, in the
release built from the CRSP database of August 2026 and retrieved on 19 October 2026; these files are
not redistributed, and \texttt{scripts/build\_french.py --download} rebuilds them. The US 10 year
yield, the WTI price and the exchange rates are FRED series \parencite{fredch9}, and the Canadian yields
are from the Bank of Canada \parencite{bocvalet}; \texttt{scripts/pull\_fred.py} downloads the FRED
series and cuts them at 30 September 2026. The bitcoin and ether prices are Coinbase series published on
FRED \parencite{coinbasefred} (\texttt{CBBTCUSD} and \texttt{CBETHUSD}); they are not redistributed with
the code, and the same script rebuilds them. The bitcoin series starts on \btSeriesStart{} and has gaps in its first weeks and no values from the day after \btStartDay{} until \btNextDay, so
its December 2014 price is that of \btStartDay; every month is sampled at its last
available daily value. The FRED market series are \texttt{DGS10}, \texttt{DCOILWTICO}, \texttt{DEXUSEU},
\texttt{DEXJPUS} and \texttt{DEXCAUS}, and the Bank of Canada series are \texttt{BD.CDN.10YR.DQ.YLD},
\texttt{BD.CDN.7YR.DQ.YLD} and \texttt{TB.CDN.90D.MID}. The carry section uses the Chapter 5 file of the USD/CAD rate and the Canadian and US one month
bill yields, built from Bank of Canada series and the FRED series \texttt{DGS1MO}, through September 2026,
and the 7 and 10 year Canadian benchmark yields. The yearly equity volatility of \cref{fig9:onchain} uses
the daily market return of the French library. The funds of
\cref{sec9:biases,sec9:replication,sec9:fees,sec9:fof}, the oil curve and the exercise inputs are
illustrations.

\printbibliography[heading=subbibliography,title={References for Chapter \thechapter}]
\end{refsection}
